% Options for packages loaded elsewhere
\PassOptionsToPackage{unicode}{hyperref}
\PassOptionsToPackage{hyphens}{url}
\PassOptionsToPackage{dvipsnames,svgnames,x11names}{xcolor}
%
\documentclass[
  11pt,
]{article}
\usepackage{amsmath,amssymb}
\usepackage{iftex}
\ifPDFTeX
  \usepackage[T1]{fontenc}
  \usepackage[utf8]{inputenc}
  \usepackage{textcomp} % provide euro and other symbols
\else % if luatex or xetex
  \usepackage{unicode-math} % this also loads fontspec
  \defaultfontfeatures{Scale=MatchLowercase}
  \defaultfontfeatures[\rmfamily]{Ligatures=TeX,Scale=1}
\fi
\usepackage{lmodern}
\ifPDFTeX\else
  % xetex/luatex font selection
\fi
% Use upquote if available, for straight quotes in verbatim environments
\IfFileExists{upquote.sty}{\usepackage{upquote}}{}
\IfFileExists{microtype.sty}{% use microtype if available
  \usepackage[]{microtype}
  \UseMicrotypeSet[protrusion]{basicmath} % disable protrusion for tt fonts
}{}
\makeatletter
\@ifundefined{KOMAClassName}{% if non-KOMA class
  \IfFileExists{parskip.sty}{%
    \usepackage{parskip}
  }{% else
    \setlength{\parindent}{0pt}
    \setlength{\parskip}{6pt plus 2pt minus 1pt}}
}{% if KOMA class
  \KOMAoptions{parskip=half}}
\makeatother
\usepackage{xcolor}
\usepackage{microtype}
\usepackage[a4paper,margin=1in]{geometry}
\usepackage{color}
\usepackage{fancyvrb}
\newcommand{\VerbBar}{|}
\newcommand{\VERB}{\Verb[commandchars=\\\{\}]}
\DefineVerbatimEnvironment{Highlighting}{Verbatim}{commandchars=\\\{\},fontsize=\small}
% Add ',fontsize=\small' for more characters per line
\newenvironment{Shaded}{}{}
\newcommand{\AlertTok}[1]{\textcolor[rgb]{1.00,0.00,0.00}{\textbf{#1}}}
\newcommand{\AnnotationTok}[1]{\textcolor[rgb]{0.38,0.63,0.69}{\textbf{\textit{#1}}}}
\newcommand{\AttributeTok}[1]{\textcolor[rgb]{0.49,0.56,0.16}{#1}}
\newcommand{\BaseNTok}[1]{\textcolor[rgb]{0.25,0.63,0.44}{#1}}
\newcommand{\BuiltInTok}[1]{\textcolor[rgb]{0.00,0.50,0.00}{#1}}
\newcommand{\CharTok}[1]{\textcolor[rgb]{0.25,0.44,0.63}{#1}}
\newcommand{\CommentTok}[1]{\textcolor[rgb]{0.38,0.63,0.69}{\textit{#1}}}
\newcommand{\CommentVarTok}[1]{\textcolor[rgb]{0.38,0.63,0.69}{\textbf{\textit{#1}}}}
\newcommand{\ConstantTok}[1]{\textcolor[rgb]{0.53,0.00,0.00}{#1}}
\newcommand{\ControlFlowTok}[1]{\textcolor[rgb]{0.00,0.44,0.13}{\textbf{#1}}}
\newcommand{\DataTypeTok}[1]{\textcolor[rgb]{0.56,0.13,0.00}{#1}}
\newcommand{\DecValTok}[1]{\textcolor[rgb]{0.25,0.63,0.44}{#1}}
\newcommand{\DocumentationTok}[1]{\textcolor[rgb]{0.73,0.13,0.13}{\textit{#1}}}
\newcommand{\ErrorTok}[1]{\textcolor[rgb]{1.00,0.00,0.00}{\textbf{#1}}}
\newcommand{\ExtensionTok}[1]{#1}
\newcommand{\FloatTok}[1]{\textcolor[rgb]{0.25,0.63,0.44}{#1}}
\newcommand{\FunctionTok}[1]{\textcolor[rgb]{0.02,0.16,0.49}{#1}}
\newcommand{\ImportTok}[1]{\textcolor[rgb]{0.00,0.50,0.00}{\textbf{#1}}}
\newcommand{\InformationTok}[1]{\textcolor[rgb]{0.38,0.63,0.69}{\textbf{\textit{#1}}}}
\newcommand{\KeywordTok}[1]{\textcolor[rgb]{0.00,0.44,0.13}{\textbf{#1}}}
\newcommand{\NormalTok}[1]{#1}
\newcommand{\OperatorTok}[1]{\textcolor[rgb]{0.40,0.40,0.40}{#1}}
\newcommand{\OtherTok}[1]{\textcolor[rgb]{0.00,0.44,0.13}{#1}}
\newcommand{\PreprocessorTok}[1]{\textcolor[rgb]{0.74,0.48,0.00}{#1}}
\newcommand{\RegionMarkerTok}[1]{#1}
\newcommand{\SpecialCharTok}[1]{\textcolor[rgb]{0.25,0.44,0.63}{#1}}
\newcommand{\SpecialStringTok}[1]{\textcolor[rgb]{0.73,0.40,0.53}{#1}}
\newcommand{\StringTok}[1]{\textcolor[rgb]{0.25,0.44,0.63}{#1}}
\newcommand{\VariableTok}[1]{\textcolor[rgb]{0.10,0.09,0.49}{#1}}
\newcommand{\VerbatimStringTok}[1]{\textcolor[rgb]{0.25,0.44,0.63}{#1}}
\newcommand{\WarningTok}[1]{\textcolor[rgb]{0.38,0.63,0.69}{\textbf{\textit{#1}}}}
\usepackage{longtable,booktabs,array}
\usepackage{calc} % for calculating minipage widths
% Correct order of tables after \paragraph or \subparagraph
\usepackage{etoolbox}
\makeatletter
\patchcmd\longtable{\par}{\if@noskipsec\mbox{}\fi\par}{}{}
\makeatother
% Allow footnotes in longtable head/foot
\IfFileExists{footnotehyper.sty}{\usepackage{footnotehyper}}{\usepackage{footnote}}
\makesavenoteenv{longtable}
\usepackage{graphicx}
\usepackage{float}
\usepackage{caption}
\makeatletter
\def\maxwidth{\ifdim\Gin@nat@width>\linewidth\linewidth\else\Gin@nat@width\fi}
\def\maxheight{\ifdim\Gin@nat@height>\textheight\textheight\else\Gin@nat@height\fi}
\makeatother
% Scale images if necessary, so that they will not overflow the page
% margins by default, and it is still possible to overwrite the defaults
% using explicit options in \includegraphics[width, height, ...]{}
\setkeys{Gin}{width=\maxwidth,height=\maxheight,keepaspectratio}
% Set default figure placement to htbp
\makeatletter
\def\fps@figure{htbp}
\makeatother
\setlength{\emergencystretch}{8em} % prevent overfull lines
\providecommand{\tightlist}{%
  \setlength{\itemsep}{0pt}\setlength{\parskip}{0pt}}
\setcounter{secnumdepth}{-\maxdimen} % remove section numbering
\ifLuaTeX
\usepackage[bidi=basic]{babel}
\else
\usepackage[bidi=default]{babel}
\fi
\babelprovide[main,import]{spanish}
% get rid of language-specific shorthands (see #6817):
\let\LanguageShortHands\languageshorthands
\def\languageshorthands#1{}
\ifLuaTeX
  \usepackage{selnolig}  % disable illegal ligatures
\fi
\usepackage{bookmark}
\IfFileExists{xurl.sty}{\usepackage{xurl}}{} % add URL line breaks if available
\urlstyle{same}
\hypersetup{
  pdftitle={Particiones de aristas en cliques de grafos split con error lineal mediante empaquetamientos estructurados de triángulos},
  pdflang={es},
  colorlinks=true,
  linkcolor={blue},
  filecolor={Maroon},
  citecolor={Blue},
  urlcolor={Blue},
  pdfcreator={LaTeX via pandoc}}

\title{Particiones de aristas en cliques de grafos split con error lineal mediante empaquetamientos estructurados de triángulos}
\author{Juan Pablo Traverso Gianini\\
\small Investigador independiente, Santiago, Chile\\
\small \href{mailto:jtraverso@gmail.com}{jtraverso@gmail.com}\\
\small ORCID: 0009-0003-6068-4096}
\date{21 de julio de 2026}

\begin{document}
\sloppy
\maketitle
\begin{small}\begin{raggedright}

\textbf{Paper III de la serie}\\
\textbf{Manuscrito de autor:} v0.9.5, borrador de revisión con estilo unificado\\
\textbf{Fecha:} 21 de julio de 2026\\
\textbf{Estado:} manuscrito de trabajo; verificación formal en curso. El proyecto Lean 4 / Mathlib v4.28.0 compila actualmente como un grafo completo de dependencias. Su núcleo elemental se ha reducido a ocho obligaciones de prueba localizadas, sin \texttt{admit}, \texttt{native\_decide}, \texttt{unsafe} ni axiomas de proyecto no registrados. Los únicos axiomas matemáticos previstos son los dos insumos asintóticos externos enunciados en la Sección 2.4. Las expresiones de liberación ``formalmente verificado'' y ``libre de \texttt{sorry}'' se reservan hasta congelar el build final, el informe de axiomas y el commit. No ha sido sometido a revisión externa; la revisión bibliográfica y de novedad está incompleta. \textbf{MSC 2020:} Primario 05C70; Secundarios 05C35, 05C72, 05C15.

\begin{center}\rule{0.5\linewidth}{0.5pt}\end{center}

\end{raggedright}\end{small}

\section{Resumen}\label{resumen}

Sea \(G=(K\cup I,E)\) un grafo split sobre \(n\) vértices, donde \(K\)
es una clique e \(I\) es un conjunto independiente. Probamos que existe
una constante absoluta \(C\) tal que

\[
|E(G)|-2\nu_3(G)
\le
\frac{n^2}{6}+Cn,
\]

donde \(\nu_3(G)\) es el número máximo de triángulos dos a dos
arista-disjuntos en \(G\). En consecuencia,

\[
\operatorname{cp}(G)
\le
\frac{n^2}{6}+Cn
\]

para todo grafo split \(G\), donde \(\operatorname{cp}(G)\) es el número
mínimo de cliques cuyos conjuntos de aristas particionan \(E(G)\). Esto
refuerza la cota asintótica \((1/6+o(1))n^2\) hasta la escala correcta
de error lineal. La constante principal es óptima.

La demostración separa tres regímenes según \(\alpha=|I|/|K|\). En el
régimen bulk, un programa lineal exacto de cuatro órbitas para un perfil
de vecindario común produce un margen fraccionario cuadrático uniforme;
el teorema de Haxell--Rödl absorbe entonces la pérdida subcuadrática de
integralidad. Cuando \(\alpha\to0\), grandes matchings arista-disjuntos
anclados en los vértices independientes dejan un residuo de clique casi
completo y triangularmente divisible, que se descompone exactamente
mediante resultados de descomposición de grafos densos. Cuando
\(\alpha\to2\), una factorización promediada cierra el corredor corto,
mientras que una desigualdad de polarización con doble factor y un
argumento de completación de ganancias respecto de un centro desplazado
cierran el corredor mesoscópico restante.

La demostración no establece la estimación universal más fuerte
\(\nu_3^*(G)-\nu_3(G)=O(n)\). Fuera del corredor extremal, una brecha de
integralidad potencialmente superlineal es absorbida por la holgura
fraccionaria cuadrática.

Todos los mecanismos del corredor son autocontenidos; en particular, el
ingrediente de coloreo de aristas por listas se demuestra desde primeros
principios en el Apéndice D. Solo quedan dos teoremas asintóticos
externos, confinados a los regímenes bulk y escaso, y el teorema del
corredor casi extremal (Proposición 10.1) es completamente elemental y
efectivo.

\textbf{Palabras clave:} partición en cliques; grafo split;
empaquetamiento de triángulos; empaquetamiento fraccionario;
factorización; polarización; coloreo de aristas por listas;
descomposición de grafos.

\begin{center}\rule{0.5\linewidth}{0.5pt}\end{center}

\section{1. Introducción}\label{introducciuxf3n}

\subsection{1.1 Particiones en cliques de grafos
split}\label{particiones-en-cliques-de-grafos-split}

Una \textbf{partición en cliques} de un grafo \(G\) es una familia de
subgrafos completos cuyos conjuntos de aristas particionan \(E(G)\). Se
permiten cliques de orden dos. El tamaño mínimo de una familia de este
tipo se denota por \(\operatorname{cp}(G)\).

El problema de Erdős, Ordman y Zalcstein sobre particiones en cliques de
grafos cordales pregunta por el orden extremal correcto de
\(\operatorname{cp}(G)\) en esa clase. Su construcción complete-split
---de hecho, threshold--- muestra que el coeficiente \(1/6\) es
inevitable. Los grafos split forman una subclase natural e importante de
los grafos cordales: un grafo es split si su conjunto de vértices puede
particionarse en una clique y un conjunto independiente.

Este es el tercer paper de la serie sobre el problema de
Erdős--Ordman--Zalcstein de particiones en cliques (Problema de Erdős
n.º 81). El Paper I prueba una cota fraccionaria finita para grafos
split, \[
|E(G)|-2\nu_3^*(G)\le n^2/6+n.
\] El Paper II identifica la familia terminal complete-split y determina
el máximo finito exacto de la funcional fraccionaria de cobertura sobre
grafos cordales. Los Papers I y II son fraccionarios y están formalmente
verificados en Lean 4. El presente paper es la primera entrega integral:
refuerza la estimación para grafos split hasta la escala de error lineal
\(n^2/6+O(n)\).

El Paper II determina exactamente el problema extremal fraccionario.
Para todo \(n\),

\[
\max_{\substack{|V(G)|=n\\G\text{ cordal}}}
\bigl(|E(G)|-2\nu_3^*(G)\bigr)
=
\max_{\substack{|V(G)|=n\\G\text{ split}}}
\bigl(|E(G)|-2\nu_3^*(G)\bigr)
=
\left\lfloor\frac{(2n+1)^2}{24}\right\rfloor.
\]

Por el redondeo fraccionario-a-integral estándar, se sigue que todo
grafo cordal satisface

\[
\operatorname{cp}(G)
\le
\left(\frac16+o(1)\right)n^2.
\]

El presente paper aborda una cuestión distinta y genuinamente de segundo
orden: refuerza la cota integral hasta \(n^2/6+O(n)\) para grafos split
arbitrarios.

\subsection{1.2 Empaquetamientos de triángulos y objetivo de error
lineal}\label{empaquetamientos-de-triuxe1ngulos-y-objetivo-de-error-lineal}

Sea \(\nu_3(G)\) el número máximo de triángulos dos a dos
arista-disjuntos en \(G\). Todo empaquetamiento de este tipo produce una
partición en cliques: se conservan los triángulos empaquetados y se usa
un \(K_2\) para cada arista no cubierta. Por tanto,

\[
\operatorname{cp}(G)
\le
|E(G)|-2\nu_3(G).
\tag{1.1}
\]

El teorema central de este paper es una estimación integral en la escala
correcta de segundo orden.

\subsubsection{Teorema 1.1}\label{teorema-1.1-teorema-de-error-lineal-para-grafos-split}

Existe una constante absoluta \(C\) tal que todo grafo split \(G\) sobre
\(n\) vértices satisface

\[
\boxed{
|E(G)|-2\nu_3(G)
\le
\frac{n^2}{6}+Cn.
}
\tag{1.2}
\]

Al combinar (1.1) con el Teorema 1.1 se obtiene la afirmación para
particiones en cliques.

\subsubsection{Corolario 1.2}\label{corolario-1.2-cota-de-error-lineal-para-particiones-en-cliques}

Existe una constante absoluta \(C\) tal que todo grafo split \(G\) sobre
\(n\) vértices satisface

\[
\boxed{
\operatorname{cp}(G)
\le
\frac{n^2}{6}+Cn.
}
\tag{1.3}
\]

La familia complete-split ---de hecho threshold--- de
Erdős--Ordman--Zalcstein \[
K_p\vee\overline K_{2p},
\qquad n=3p,
\] muestra que la constante principal \(1/6\) es óptima y que aparece de
manera natural un término lineal.

\begin{figure}[H]
\centering
\includegraphics[width=0.72\textwidth]{paperIII_figures_bilingual_v0.9.7/fig1_complete_split_extremizer_es.png}
\caption{El grafo complete-split \(K_p\vee\overline K_{2p}\). Todas las aristas entre las dos partes están presentes y el lado independiente no tiene aristas internas. Como todo triángulo contiene una arista de la clique, un empaquetamiento de triángulos arista-disjuntos utiliza a lo sumo \(\binom p2\) triángulos; véase la Sección 10.2. La figura es ilustrativa y no se usa como premisa de la demostración.}
\label{fig:complete-split}
\end{figure}

\subsection{1.3 Localización de obstrucciones
hipotéticas}\label{localizaciuxf3n-de-obstrucciones-hipotuxe9ticas}

\subsubsection{Proposición 1.3}\label{proposiciuxf3n-1.3-localizaciuxf3n-de-obstrucciones-hipotuxe9ticas}

Supongamos, por contradicción, que no existe una cota absoluta con error
lineal y, para cada entero positivo \(k\), escojamos un grafo split
\(G_k\) de orden mínimo que satisface

\[
|E(G_k)|-2\nu_3(G_k)>\frac{|V(G_k)|^2}{6}+k|V(G_k)|.
\]

Tras pasar a una subsucesión y escribir \(|K_k|=p_k\), \(|I_k|=q_k\),
toda sucesión hipotética de este tipo debe satisfacer

\[
\frac{q_k}{p_k}\longrightarrow 2.
\]

Además, si \(q_k=2p_k-s_k\), las estimaciones del régimen bulk y del
corredor corto fuerzan

\[
\sqrt{p_k}\ll s_k=o(p_k).
\]

Por tanto, toda obstrucción a un término de error lineal tendría que
concentrarse en el corredor mesoscópico casi extremal. La dicotomía
entre dispersión alta y baja de la Sección 9 excluye ese corredor.

\subsection{1.4 Lo que el teorema no
demuestra}\label{lo-que-el-teorema-no-demuestra}

Sea \(\nu_3^*(G)\) el número fraccionario de empaquetamiento de
triángulos. La presente demostración \textbf{no} establece

\[
\nu_3^*(G)-\nu_3(G)=O(n)
\tag{1.4}
\]

uniformemente sobre todos los grafos split.

En efecto, escribamos

\[
\Delta(G)=\nu_3^*(G)-\nu_3(G)
\]

y

\[
S(G)=\frac{n^2}{6}-\bigl(|E(G)|-2\nu_3^*(G)\bigr).
\]

Entonces

\[
|E(G)|-2\nu_3(G)
=
\frac{n^2}{6}-S(G)+2\Delta(G).
\tag{1.5}
\]

El Teorema 1.1 solo requiere

\[
2\Delta(G)\le S(G)+O(n).
\]

En el régimen bulk, \(S(G)\) es cuadrático y absorbe la pérdida general
de integralidad \(o(n^2)\). Por ello, el problema de una brecha
universal lineal de integralidad permanece abierto.

\subsection{1.5 Arquitectura de la
demostración}\label{arquitectura-de-la-demostraciuxf3n}

Escribamos

\[
|K|=p,
\qquad
|I|=q,
\qquad
\alpha=\frac qp.
\]

La razón \(\alpha=q/p\) determina qué fuente de holgura está disponible.
Primero resolvemos \(q\ge2p-1\) mediante una factorización promediada
directa, y desde entonces podemos suponer \(0\le\alpha<2\). Una sucesión
hipotética de contraejemplos tiene una subsucesión en uno de tres
regímenes.

\begin{enumerate}
\def\labelenumi{\arabic{enumi}.}
\item
  \textbf{Bulk:} \(\alpha\) permanece alejado tanto de \(0\) como de
  \(2\). Un LP exacto de perfil común y un argumento de clonado
  fraccional producen holgura fraccionaria cuadrática. Haxell--Rödl
  convierte el empaquetamiento fraccionario en uno integral con pérdida
  \(o(n^2)\).
\item
  \textbf{Lado independiente escaso:} \(\alpha\to0\). Cada vértice
  independiente ancla un matching grande dentro de su vecindario. La
  clique residual es casi completa. Tras eliminar \(O(p)\) aristas para
  corregir divisibilidad, admite una descomposición exacta en
  triángulos.
\item
  \textbf{Corredor casi extremal:} \(\alpha\to2\). Escribimos
  \(q=2p-s\). La factorización promediada cierra \(s=O(\sqrt p)\). Para
  \(\sqrt p\ll s=o(p)\), la dispersión alta se paga mediante un término
  de polarización con doble factor, mientras que la dispersión baja
  implica cercanía a un centro común y se trata mediante completación de
  ganancias respecto de un centro desplazado.
\end{enumerate}

La división no es meramente formal. En el bulk hay suficiente holgura
fraccionaria cuadrática para absorber una pérdida general de redondeo
\(o(n^2)\). En los dos extremos esa holgura desaparece, y la
demostración reemplaza el redondeo genérico por construcciones adaptadas
a la geometría de la presentación split.

\begin{figure}[H]
\centering
\includegraphics[width=0.94\textwidth]{paperIII_figures_bilingual_v0.9.7/fig2_alpha_regimes_es.png}
\caption{Los tres regímenes usados en la demostración. Se muestran los umbrales principales del corredor; las hipótesis completas aparecen en las Secciones 4, 8 y 9. Los regímenes extremos son asintóticos, no subintervalos rígidos de \([0,2]\).}
\label{fig:alpha-regimes}
\end{figure}

\subsection{1.6 Auditorías computacionales
suplementarias}\label{auditoruxedas-computacionales-suplementarias}

El paquete de cierre adjunto contiene pruebas de regresión con
aritmética exacta e instancias finitas para identidades y desigualdades
cuantitativas seleccionadas de la demostración.

Cada afirmación auditada se demuestra analíticamente en el manuscrito.
Ningún cálculo, script, enumeración finita ni salida de solver es una
premisa lógica del Teorema 1.1. El material suplementario se proporciona
únicamente para verificación independiente, pruebas de regresión y
reproducibilidad.

\subsection{1.7 Organización}\label{organizaciuxf3n}

La Sección 2 fija la notación y los insumos externos. La Sección 3
resuelve el LP de perfil común. La Sección 4 demuestra la cota de
clonado fraccional y el margen fraccionario global. La Sección 5
desarrolla el redondeo por factorización simple y doble. La Sección 6
demuestra la polarización. La Sección 7 demuestra la completación de
ganancias respecto de un centro desplazado. La Sección 8 trata el
régimen con lado independiente escaso. La Sección 9 ensambla los tres
regímenes. La Sección 10 registra los corolarios, la Sección 11 discute
la demostración y sus limitaciones, y la Sección 12 describe usos
futuros y direcciones abiertas. Los apéndices contienen detalles
algebraicos, la corrección de divisibilidad, material de
reproducibilidad y una demostración autocontenida del teorema de coloreo
de aristas por listas (Apéndice D).

\begin{center}\rule{0.5\linewidth}{0.5pt}\end{center}

\section{2. Preliminares}\label{preliminares}

\subsection{2.1 Notación split}\label{notaciuxf3n-split}

En todo el paper,

\[
G=(K\sqcup I,E)
\]

es un grafo split, donde \(K\) es una clique de tamaño \(p\) e
\(I=\{v_1,\ldots,v_q\}\) es independiente.

Para \(v_i\in I\), escribimos

\[
N_i=N(v_i)\cap K,
\qquad
d_i=|N_i|,
\qquad
S_i=K\setminus N_i,
\qquad
m_i=|S_i|.
\]

Definimos

\[
M=\sum_i m_i,
\qquad
S_2=\sum_i m_i^2.
\]

Cuando \(q\) está cerca de \(2p\), escribimos

\[
q=2p-s.
\tag{2.1}
\]

\subsection{2.2 Empaquetamientos de
triángulos}\label{empaquetamientos-de-triuxe1ngulos}

Un empaquetamiento fraccionario de triángulos es una asignación de pesos
no negativos a los triángulos de \(G\) tal que el peso total que pasa
por cada arista es a lo sumo uno. Su valor máximo es \(\nu_3^*(G)\).

El dual es una cobertura fraccionaria de triángulos: una asignación de
pesos no negativos a las aristas tal que cada triángulo recibe peso
total al menos uno. La dualidad de programación lineal da igualdad entre
ambos valores óptimos.

\subsection{2.3 Factorizaciones de grafos
completos}\label{factorizaciones-de-grafos-completos}

Adoptamos las convenciones de borde

\[
\chi'(K_0)=\chi'(K_1)=0.
\]

Para \(t\ge2\), el índice cromático del grafo completo es

\[
\chi'(K_t)
=
\begin{cases}
t-1,&t\text{ par},\\
t,&t\text{ impar}.
\end{cases}
\tag{2.2}
\]

Por tanto, \(E(K_t)\) se descompone en \(\chi'(K_t)\) matchings. Las
convenciones para \(t\le1\) hacen que el mismo lenguaje sea válido para
conjuntos vacíos de aristas.

\subsection{2.4 Teoremas externos}\label{teoremas-externos}

Usamos dos insumos asintóticos externos. Versiones anteriores citaban un
tercero, el teorema de coloreo de aristas por listas; el caso requerido
por este paper se demuestra de manera autocontenida en el Apéndice D,
por lo que las Secciones 5--7 y la Proposición 10.1 no dependen de
teoremas externos.

\subsubsection{Teorema 2.1 ---
Haxell--Rödl/Yuster}\label{teorema-2.1-haxellruxf6dlyuster}

Para todo grafo fijo \(H\),

\[
\nu_H^*(G)-\nu_H(G)=o(|V(G)|^2)
\]

uniformemente sobre los grafos \(G\). Lo aplicamos con \(H=K_3\).

\subsubsection{Teorema 2.2}\label{teorema-2.2-coloreo-de-aristas-por-listas-demostrado-en-el-apuxe9ndice-d}

Sea \(B\) un grafo bipartito de grado máximo \(\Delta(B)\). Si a cada
arista \(e\) se le asigna una lista \(L(e)\) con

\[
|L(e)|\ge\Delta(B),
\]

entonces \(B\) admite un coloreo propio de aristas que elige el color de
cada arista desde su lista.

Este es el caso de grado máximo del teorema de Galvin {[}12{]}. En el
Apéndice D se entrega una demostración autocontenida ---coloreo de
König, emparejamientos estables y lema del núcleo---, por lo que no
constituye una dependencia externa del paper. No se necesita el
refinamiento local \[
|L(xy)|\ge\max\{d(x),d(y)\}
\] de Borodin, Kostochka y Woodall {[}2{]}: la hipótesis (7.2) ya acota
cada lista por el grado máximo del grafo de ganancias.

\subsubsection{Teorema 2.3}\label{teorema-2.3-descomposiciuxf3n-densa-en-triuxe1ngulos}

Para todo \(\varepsilon>0\), todo grafo triangularmente divisible \(H\),
suficientemente grande, que satisface

\[
\delta(H)\ge(0.9+\varepsilon)|V(H)|
\]

admite una descomposición en triángulos. Esto se sigue del teorema de
Dross sobre descomposición fraccionaria en triángulos ---grado mínimo
\(0.9v\) es suficiente fraccionariamente--- combinado con el teorema de
absorción iterativa de Barber, Kühn, Lo y Osthus, que convierte el
umbral fraccionario en un umbral de descomposición exacta para grafos
triangularmente divisibles, salvo \(\varepsilon\) y para orden
suficientemente grande.

\begin{center}\rule{0.5\linewidth}{0.5pt}\end{center}

\section{3. El programa lineal de perfil
común}\label{el-programa-lineal-de-perfil-comuxfan}

Para enteros \(p,q,d\), sea \(H(p,q,d)\) el grafo split con clique
\(K\), \(|K|=p\), conjunto independiente \(I\), \(|I|=q\), y un conjunto
fijo \(N\subseteq K\), \(|N|=d\), tal que todo vértice de \(I\) tiene
vecindario \(N\). Definimos \(R=K\setminus N\) y \(r=p-d\).

\subsection{3.1 Variables simétricas de
cobertura}\label{variables-simuxe9tricas-de-cobertura}

Al promediar sobre permutaciones de \(N\), \(R\) e \(I\), podemos
suponer que una cobertura fraccionaria óptima de triángulos es constante
en las cuatro clases de aristas

\[
E(N),
\qquad E(N,I),
\qquad E(N,R),
\qquad E(R).
\]

Sean \(a,b,c,e\) sus respectivos pesos. Las restricciones de triángulos
son

\[
a+2b\ge1,
\qquad
3a\ge1,
\tag{3.1}
\]

\[
a+2c\ge1,
\qquad
2c+e\ge1,
\qquad
3e\ge1.
\tag{3.2}
\]

Las restricciones correspondientes a tipos de triángulos vacíos pueden
omitirse; la fórmula siguiente sigue siendo válida para \(p\ge3\).

La función objetivo es

\[
\binom d2a+qdb+dr\,c+\binom r2e.
\tag{3.3}
\]

\subsection{3.2 Solución exacta}\label{soluciuxf3n-exacta}

\subsubsection{Teorema 3.1}\label{teorema-3.1-fuxf3rmula-de-perfil-comuxfan}

Para \(p\ge3\),

\[
\boxed{
\nu_3^*(H(p,q,d))
=
F(p,q,d),
}
\tag{3.4}
\]

donde

\[
\boxed{
F(p,q,d)
=
\min\left\{
\frac{\binom p2+qd}{3},
\binom d2+\binom r2,
\binom d2+\frac{dr+\binom r2}{3}
\right\}.
}
\tag{3.5}
\]

\subsubsection{Demostración}\label{demostraciuxf3n}

En un óptimo,

\[
b=\frac{1-a}{2},
\qquad
c=\frac{1-\min\{a,e\}}2,
\qquad
\frac13\le a,e\le1.
\]

Si \(a\le e\), disminuir \(e\) hasta \(a\) no puede aumentar la función
objetivo. Por tanto, basta minimizar sobre

\[
\frac13\le e\le a\le1.
\]

La función objetivo es afín en este triángulo, por lo que un mínimo se
alcanza en

\[
(a,e)=\left(\frac13,\frac13\right),
\qquad
(1,1),
\qquad
\left(1,\frac13\right).
\]

Los valores correspondientes son precisamente las tres expresiones de
(3.5). La dualidad completa la demostración. \(\square\)

\subsection{3.3 Interpretación de las tres
coberturas}\label{interpretaciuxf3n-de-las-tres-coberturas}

El mínimo de (3.5) se recuerda más fácilmente mediante los tres vértices
del programa reducido de dos variables. Cada vértice representa una
forma distinta de pagar las restricciones de triángulos.

\begin{longtable}[]{@{}
  >{\raggedright\arraybackslash}p{(\columnwidth - 6\tabcolsep) * \real{0.2143}}
  >{\raggedleft\arraybackslash}p{(\columnwidth - 6\tabcolsep) * \real{0.2857}}
  >{\raggedleft\arraybackslash}p{(\columnwidth - 6\tabcolsep) * \real{0.2857}}
  >{\raggedright\arraybackslash}p{(\columnwidth - 6\tabcolsep) * \real{0.2143}}@{}}
\toprule\noalign{}
\begin{minipage}[b]{\linewidth}\raggedright
Patrón de cobertura
\end{minipage} & \begin{minipage}[b]{\linewidth}\raggedleft
Vértice reducido \((a,e)\)
\end{minipage} & \begin{minipage}[b]{\linewidth}\raggedleft
Valor
\end{minipage} & \begin{minipage}[b]{\linewidth}\raggedright
Descripción geométrica
\end{minipage} \\
\midrule\noalign{}
\endhead
\bottomrule\noalign{}
\endlastfoot
\textbf{Uniforme} & \((1/3,1/3)\) & \(\bigl(\binom p2+qd\bigr)/3\) &
Todas las órbitas de aristas relevantes se cubren a la tasa fraccionaria
uniforme. \\
\textbf{Separada} & \((1,1)\) & \(\binom d2+\binom r2\) & Los dos
bloques de clique \(N\) y \(R\) se pagan internamente, sin peso en la
órbita cruzada. \\
\textbf{Vecindario caliente} & \((1,1/3)\) &
\(\binom d2+\bigl(dr+\binom r2\bigr)/3\) & El vecindario común \(N\) se
paga por completo, mientras la estructura residual de la clique
permanece fraccionaria. \\
\end{longtable}

Las etiquetas son descriptivas y no introducen nuevas definiciones. En
los cálculos posteriores nos referimos a las ramas primera, segunda y
tercera de \(F\). Esta tricotomía finita es la fuente del margen
fraccionario global.

\begin{center}\rule{0.5\linewidth}{0.5pt}\end{center}

\section{4. Clonado fraccional y margen fraccionario
exacto}\label{clonado-fraccional-y-margen-fraccionario-exacto}

\subsection{4.1 Clonado fraccional}\label{clonado-fraccional}

La operación de clonado usada aquí es el análogo, en el lado de
cobertura, de la simetrización por clones del Paper II. Allí se copian
por pares vértices o clases de clones para simplificar un grafo cordal.
Aquí se copia \(q\) veces un perfil del lado independiente en un único
paso de promedio; el grafo \(H(p,q,d_i)\) registra el perfil común
resultante.

\subsubsection{Lema 4.1}\label{lema-4.1-cota-de-clonado-fraccional}

Para todo grafo split con \(q\ge1\),

\[
\boxed{
\nu_3^*(G)
\ge
\frac1q\sum_{i=1}^qF(p,q,d_i).
}
\tag{4.1}
\]

\subsubsection{Demostración}\label{demostraciuxf3n-1}

Sea \(y\) cualquier cobertura fraccionaria de triángulos de \(G\).
Definamos

\[
A=\sum_{e\in E(K)}y_e
\]

y

\[
B_i=\sum_{x\in N_i}y_{v_ix}.
\]

Reemplazamos \(v_i\) por \(q\) clones independientes, todos con
vecindario \(N_i\), y damos a cada clon los pesos incidentes de \(v_i\).
Junto con los pesos originales de las aristas de la clique, esto
constituye una cobertura fraccionaria de \(H(p,q,d_i)\) de peso
\(A+qB_i\). Por tanto,

\[
A+qB_i\ge F(p,q,d_i).
\]

Al sumar sobre \(i\), obtenemos

\[
q\left(A+\sum_iB_i\right)
\ge
\sum_iF(p,q,d_i).
\]

Minimizar sobre las coberturas demuestra el lema. \(\square\)

\subsection{4.2 El margen exacto}\label{el-margen-exacto}

Supongamos \(0<q\le2p\), definamos

\[
\alpha=\frac qp,
\]

y el umbral de empaquetamiento

\[
T(G)=\frac12\left(|E(G)|-\frac{(p+q)^2}{6}\right).
\tag{4.2}
\]

Definamos

\[
\mu(\alpha)
=
\begin{cases}
\alpha^2/12,&0\le\alpha\le2/3,\\
(2-\alpha)^2/48,&2/3\le\alpha\le2.
\end{cases}
\tag{4.3}
\]

\subsubsection{Teorema 4.2}\label{teorema-4.2-margen-fraccionario-unificado}

\[
\boxed{
\nu_3^*(G)
\ge
T(G)+\mu(\alpha)p^2-\frac p4.
}
\tag{4.4}
\]

\subsubsection{Demostración}\label{demostraciuxf3n-2}

Sea

\[
C_\alpha=\frac{2-2\alpha-\alpha^2}{12}.
\]

Para cada rama de \(F(p,q,d)\), completar cuadrados da

\[
F(p,q,d)
\ge
\frac{qd}{2}+C_\alpha p^2+\mu(\alpha)p^2-\frac p2.
\tag{4.5}
\]

Los tres mínimos residuales normalizados son

\[
\frac{\alpha^2}{12},
\qquad
\frac{(2-\alpha)^2}{48},
\]

y

\[
\begin{cases}
\alpha(8-5\alpha)/48,&\alpha\le4/3,\\
(2-\alpha)^2/12,&\alpha\ge4/3.
\end{cases}
\]

El tercero nunca está por debajo del mínimo de los dos primeros. Al
promediar (4.5) mediante el Lema 4.1 obtenemos

\[
\nu_3^*(G)
\ge
\frac12\sum_i d_i+C_\alpha p^2+\mu(\alpha)p^2-\frac p2.
\]

Como

\[
T(G)=\frac12\sum_i d_i+C_\alpha p^2-\frac p4,
\]

se sigue (4.4). \(\square\)

\subsection{4.3 Consecuencia en el bulk}\label{consecuencia-en-el-bulk}

Si

\[
\varepsilon\le\alpha\le2-\varepsilon,
\]

entonces \(\mu(\alpha)\ge c_\varepsilon>0\). El Teorema 4.2 y
Haxell--Rödl implican

\[
\nu_3(G)\ge T(G)
\]

para todos los grafos suficientemente grandes de este régimen. Por
consiguiente,

\[
|E(G)|-2\nu_3(G)
\le
\frac{n^2}{6}.
\tag{4.6}
\]

\begin{center}\rule{0.5\linewidth}{0.5pt}\end{center}

\section{\texorpdfstring{5. Redondeo por factorización cerca de
\(q=2p\)}{5. Redondeo por factorización cerca de q=2p}}\label{redondeo-por-factorizaciuxf3n-cerca-de-q2p}

Sea

\[
r_p=\chi'(K_p).
\]

\subsection{5.1 Promedio con un factor}\label{promedio-con-un-factor}

\subsubsection{Lema 5.1}\label{lema-5.1}

Si \(q\ge r_p\), entonces

\[
\boxed{
\nu_3(G)
\ge
\frac1q\sum_i\binom{d_i}{2}.
}
\tag{5.1}
\]

\subsubsection{Demostración}\label{demostraciuxf3n-3}

Factorizamos \(K_p\) en \(r_p\) matchings. Asignamos los factores
inyectiva y uniformemente a los vértices de \(I\). En el factor asignado
a \(v_i\), conservamos solo las aristas con ambos extremos en \(N_i\).
Las aristas conservadas forman triángulos \(KKI\) válidos y
arista-disjuntos.

El número esperado de aristas conservadas es

\[
\frac1q\sum_i\sum_{j=1}^{r_p}|F_j\cap E(K[N_i])|
=
\frac1q\sum_i\binom{d_i}{2}.
\]

Alguna asignación alcanza al menos la esperanza. \(\square\)

Al escribir \(q=2p-s\), la sustitución da

\[
\Phi(G)
\le
\frac{n^2}{6}+\frac p2-\frac{s^2}{6}
+
\frac{(s-1)M-S_2}{q}.
\tag{5.2}
\]

Usando \(S_2\ge M^2/q\) y maximizando la parábola resultante se obtiene

\[
\boxed{
\Phi(G)
\le
\frac{n^2}{6}+\frac p2+\frac{s^2-6s+3}{12}.
}
\tag{5.3}
\]

Así, \(s=O(\sqrt p)\) queda cerrado con error lineal.

\subsection{5.2 Redondeo con doble
factor}\label{redondeo-con-doble-factor}

Para una arista \(e=xy\) de la clique, sea

\[
b_e=|\{i:\{x,y\}\not\subseteq N_i\}|.
\]

Definamos

\[
h=\min\{r_p,q-r_p\},
\qquad
\delta=\frac h{r_p}.
\]

Elegimos uniformemente al azar el conjunto de \(h\) factores que reciben
dos vértices independientes distintos, damos un vértice independiente a
cada factor restante y asignamos los vértices independientes a las
posiciones resultantes de manera uniforme e inyectiva. Todas las
esperanzas siguientes se toman respecto de ambas elecciones aleatorias;
alguna elección determinista alcanza al menos la ganancia esperada.

\subsubsection{Lema 5.2}\label{lema-5.2-desigualdad-de-doble-factor}

Si \(q\ge r_p\), entonces

\[
\boxed{
\begin{aligned}
\Phi(G)
\le{}&
\frac{n^2}{6}+\frac p2-\frac{s^2}{6}
+
\frac{(s-1)M-S_2}{q}\\
&-
\frac{2\delta V}{q(q-1)},
\end{aligned}
}
\tag{5.4}
\]

donde

\[
V=\sum_{e\in E(K)}b_e(q-b_e).
\tag{5.5}
\]

\subsubsection{Demostración}\label{demostraciuxf3n-4}

Si el factor que contiene \(e\) recibe un vértice, \(e\) se pierde con
probabilidad \(b_e/q\). Si recibe dos vértices, se pierde solo cuando
ambos son malos, con probabilidad

\[
\frac{b_e(b_e-1)}{q(q-1)}.
\]

Por tanto, el número esperado \(U\) de aristas perdidas de la clique es

\[
U
=
\frac1q\sum_e b_e
-
\frac{\delta}{q(q-1)}\sum_e b_e(q-b_e).
\]

Como cada factor es un matching, cada arista conservada puede asignarse
a uno de sus vértices admisibles sin crear una colisión de aristas. La
sustitución da (5.4). \(\square\)

\begin{center}\rule{0.5\linewidth}{0.5pt}\end{center}

\section{6. Polarización
cuantitativa}\label{polarizaciuxf3n-cuantitativa}

Para cada \(i\), definimos

\[
\mathcal B_i=\{e\in E(K):e\cap S_i\ne\varnothing\}.
\]

Entonces

\[
V=\sum_{i,j}|\mathcal B_i\setminus\mathcal B_j|.
\tag{6.1}
\]

Sea

\[
m=\max_i|S_i|.
\]

\subsubsection{Lema 6.1}\label{lema-6.1-desigualdad-de-polarizaciuxf3n}

Si \(2p-3m-1\ge0\), entonces

\[
\boxed{
V
\ge
\frac{2p-3m-1}{4}
\sum_{i,j}|S_i\triangle S_j|.
}
\tag{6.2}
\]

\subsubsection{Demostración}\label{demostraciuxf3n-5}

Sea

\[
a_{ij}=|S_i\setminus S_j|.
\]

Las aristas de \(\mathcal B_i\setminus\mathcal B_j\) son precisamente
las aristas de \(K_p-S_j\) que inciden en \(S_i\setminus S_j\). Por
tanto,

\[
|\mathcal B_i\setminus\mathcal B_j|
=
\frac{a_{ij}\bigl(2(p-|S_j|)-a_{ij}-1\bigr)}2.
\]

Como \(|S_j|\le m\) y \(a_{ij}\le m\),

\[
|\mathcal B_i\setminus\mathcal B_j|
\ge
\frac{2p-3m-1}{2}|S_i\setminus S_j|.
\]

Sumamos sobre pares ordenados y usamos

\[
\sum_{i,j}|S_i\setminus S_j|
=
\frac12\sum_{i,j}|S_i\triangle S_j|.
\]

Esto demuestra (6.2). \(\square\)

\begin{center}\rule{0.5\linewidth}{0.5pt}\end{center}

\section{7. Completación de ganancias respecto de un centro
desplazado}\label{completaciuxf3n-de-ganancias-respecto-de-un-centro-desplazado}

Fijemos \(R\subseteq K\). Definamos

\[
\rho=|R|,
\qquad
Q=K\setminus R,
\qquad
b=|Q|.
\]

Para cada \(i\), definimos

\[
T_i=S_i\setminus R,
\qquad
t_i=|T_i|,
\]

y

\[
G_i=R\setminus S_i,
\qquad
g_i=|G_i|.
\]

Sea

\[
A_R=\sum_i t_i,
\qquad
A_{2,R}=\sum_i t_i^2,
\qquad
B_R=\sum_i g_i.
\]

Definimos

\[
r_b=\chi'(K_b),
\qquad
u=q-r_b.
\]

Supongamos

\[
b\ge2,
\qquad
q\ge r_b,
\qquad
b\ge\chi'(K_\rho),
\tag{7.1}
\]

y

\[
b-t_i\ge\max\{\rho,u\}
\qquad
\text{para todo }i.
\tag{7.2}
\]

Definimos

\[
\theta_R=\frac{\max\{\rho-1,0\}}{b}
\]

y

\[
\kappa_R
=
1-2(1-\theta_R)\frac{u}{q}.
\tag{7.3}
\]

\subsection{\texorpdfstring{7.1 Empaquetamiento dentro de
\(Q\)}{7.1 Empaquetamiento dentro de Q}}\label{empaquetamiento-dentro-de-q}

Reservamos un conjunto \(U\subseteq I\) de \(u\) vértices; obsérvese que
\(|I\setminus U|=q-u=r_b\) exactamente. Asignamos biyectivamente los
\(r_b\) vértices restantes a una factorización de \(K[Q]\).

Para \(v_i\), el número de aristas no disponibles de \(K[Q]\) es

\[
\beta_i
=
\binom b2-\binom{b-t_i}{2}.
\tag{7.4}
\]

Para \(U\) fijo, el promedio sobre las asignaciones da al menos

\[
\binom b2-\frac1{r_b}\sum_{i\notin U}\beta_i
\tag{7.5}
\]

triángulos \(QQI\) arista-disjuntos.

\subsection{7.2 El grafo de ganancias}\label{el-grafo-de-ganancias}

Construimos un grafo bipartito con partes \(U\) y \(R\), uniendo \(v_i\)
con \(r\) cuando \(r\in G_i\). Asignamos la lista

\[
L(v_ir)=N_i\cap Q.
\]

Por (7.2), toda lista satisface

\[
|L(v_ir)|=b-t_i
\ge
\max\{\rho,u\}
\ge
\Delta,
\]

donde \(\Delta\) es el grado máximo del grafo de ganancias, pues
\(d(v_i)=g_i\le\rho\) y \(d(r)\le|U|=u\). El Teorema 2.2, demostrado de
manera autocontenida en el Apéndice D, proporciona un coloreo propio de
aristas por listas. Si \(v_ir\) recibe el color \(z\in Q\), tomamos el
triángulo \(v_irz\). Esto produce

\[
B_U=\sum_{i\in U}g_i
\]

triángulos \(IRQ\) arista-disjuntos.

\subsection{\texorpdfstring{7.3 Completación dentro de
\(R\)}{7.3 Completación dentro de R}}\label{completaciuxf3n-dentro-de-r}

Para cada \(z\in Q\), sea \(U_z\subseteq R\) el conjunto de vértices
\(r\) para los cuales \(rz\) fue usado por un triángulo \(IRQ\).

Si \(\rho\le1\), el grafo \(K[R]\) no tiene aristas, por lo que esta
completación no aporta triángulos \(RRQ\). Supongamos entonces
\(\rho\ge2\). Factorizamos \(K[R]\), inyectamos sus factores en los
colores \(z\in Q\) y eliminamos del factor asignado a \(z\) toda arista
incidente en \(U_z\). El promedio sobre las inyecciones pierde a lo sumo

\[
\frac{\rho-1}{b}\sum_z|U_z|
=
\theta_RB_U
\]

aristas. Por consiguiente, en todos los casos obtenemos al menos

\[
\binom\rho2-\theta_RB_U
\]

triángulos \(RRQ\) adicionales.

Las tres familias \(QQI\), \(IRQ\) y \(RRQ\) son arista-disjuntas. En
particular, la eliminación por colores prohibidos evita que un triángulo
\(IRQ\) y uno \(RRQ\) compartan una arista \(rz\).

\subsection{7.4 La desigualdad centrada}\label{la-desigualdad-centrada}

Elegimos \(U\), \(|U|=u\), que maximice

\[
\sum_{i\in U}
\left(
\frac{2\beta_i}{r_b}+2(1-\theta_R)g_i
\right).
\]

La suma de los mejores \(u\) términos es al menos \(u/q\) veces la suma
total. Usando

\[
2\sum_i\beta_i
=(2b-1)A_R-A_{2,R},
\]

obtenemos el lema local principal.

\subsubsection{Lema 7.1}\label{lema-7.1-desigualdad-de-ganancias-reservadas-respecto-de-un-centro-desplazado}

Bajo (7.1)--(7.2),

\[
\boxed{
\begin{aligned}
\Phi(G)
\le{}&
\frac{n^2}{6}+\frac p2-\frac{s^2}{6}
+s\rho-2\rho^2\\
&+
\kappa_RB_R
+
\frac{(s-2\rho-1)A_R-A_{2,R}}q.
\end{aligned}
}
\tag{7.6}
\]

\begin{center}\rule{0.5\linewidth}{0.5pt}\end{center}

\section{8. Régimen con lado independiente
escaso}\label{ruxe9gimen-con-lado-independiente-escaso}

Supongamos \(q=o(p)\). Demostramos la cota requerida de manera directa e
integral.

\subsection{8.1 Cota de grado para un contraejemplo
mínimo}\label{cota-de-grado-para-un-contraejemplo-muxednimo}

En el argumento global por contradicción, un contraejemplo mínimo con
penalización \(kn\) satisface, para todo \(v\in I\),

\[
d(v)>
\frac{2n-1}{6}+k.
\tag{8.1}
\]

Como \(q=o(p)\), esto implica eventualmente

\[
d(v)\ge2q+2.
\tag{8.2}
\]

\subsection{8.2 Matchings sucesivos}\label{matchings-sucesivos}

Ordenamos \(I=\{v_1,\ldots,v_q\}\). Elegimos sucesivamente matchings
arista-disjuntos

\[
M_i\subseteq K[N_i]
\]

de tamaño \(\lfloor d_i/2\rfloor\).

Antes de elegir \(M_i\), cada vértice ha perdido a lo sumo \(i-1\)
aristas incidentes, por lo que el grafo disponible sobre \(N_i\) tiene
grado mínimo al menos

\[
d_i-i\ge d_i/2.
\]

El teorema de Dirac proporciona un ciclo hamiltoniano y, por tanto, un
matching del tamaño requerido.

Sea

\[
F=\bigcup_iM_i.
\]

Entonces

\[
|F|
\ge
\frac12\sum_i d_i-\frac q2,
\qquad
\Delta(F)\le q.
\tag{8.3}
\]

Las aristas de \(M_i\) producen \(|M_i|\) triángulos \(KKI\)
arista-disjuntos con ápice \(v_i\).

\subsection{8.3 Corrección de
divisibilidad}\label{correcciuxf3n-de-divisibilidad}

El grafo residual sobre la clique

\[
R_0=K_p-F
\]

satisface

\[
\delta(R_0)
\ge
p-1-q
=(1-o(1))p.
\tag{8.4}
\]

En particular, para todos los miembros suficientemente grandes de la
sucesión, \(\delta(R_0)>p/2\). El teorema de Dirac da entonces un ciclo
hamiltoniano y, por tanto, un camino hamiltoniano

\[
P=x_1x_2\cdots x_p
\]

contenido en \(R_0\). Sea \(O\) el conjunto de vértices de grado impar
de \(R_0\). Por el lema del apretón de manos, \(|O|\) es par. Definimos
el subgrafo del camino \(J\subseteq P\) mediante

\[
x_jx_{j+1}\in E(J)
\quad\Longleftrightarrow\quad
|O\cap\{x_1,\ldots,x_j\}|\text{ es impar}.
\tag{8.5}
\]

Para un vértice interno \(x_j\), la paridad de su grado en \(J\) es el
cambio de paridad del prefijo entre las posiciones \(j-1\) y \(j\), y
por tanto es \(1\) exactamente cuando \(x_j\in O\). La misma conclusión
vale en los dos extremos porque \(|O|\) es par. En consecuencia,

\[
\operatorname{Odd}(J)=O,
\qquad
|E(J)|\le p-1,
\qquad
\Delta(J)\le2.
\tag{8.6}
\]

Así,

\[
R_1:=R_0-J
\]

tiene todos sus grados pares, y

\[
\delta(R_1)\ge p-1-q-2.
\tag{8.7}
\]

Como \(q=o(p)\), eventualmente \(\delta(R_1)>3p/4\). Por el teorema de
Turán, \(R_1\) contiene una copia de \(K_5\). Elegimos un \(C_4\) y un
\(C_5\) dentro de este \(K_5\) fijo. Según

\[
|E(R_1)|\equiv0,1,2\pmod3,
\]

eliminamos, respectivamente, nada, el \(C_4\) o el \(C_5\). Sea \(C\) el
ciclo eliminado, permitiendo \(C=\varnothing\), y definamos

\[
H:=R_1-E(C).
\]

Cada vértice pierde cero o dos aristas incidentes cuando se elimina un
ciclo, por lo que todos los grados de \(H\) siguen siendo pares. Además,
\(|E(C)|\equiv |E(R_1)|\pmod3\), y por tanto

\[
|E(H)|\equiv0\pmod3.
\]

Luego \(H\) es triangularmente divisible. La pérdida total de grado al
pasar de \(R_0\) a \(H\) es a lo sumo cuatro, de modo que

\[
\delta(H)\ge p-1-q-4.
\tag{8.8}
\]

Para aplicar el Teorema 2.3 con un parámetro fijo, tomemos por ejemplo
\(\varepsilon_0=1/100\). Como \(q=o(p)\), eventualmente \(q\le p/20\), y
entonces, para \(p\) suficientemente grande,

\[
\delta(H)\ge p-1-q-4\ge0.91p=(0.9+\varepsilon_0)|V(H)|.
\tag{8.9}
\]

El Teorema 2.3 da ahora una descomposición exacta de \(H\) en
triángulos. Obsérvese que el teorema de descomposición se aplica sobre
el conjunto original de \(p\) vértices: no se elimina ningún vértice
durante la corrección. La corrección completa elimina a lo sumo \(p+4\)
aristas.

\subsection{8.4 Tamaño del
empaquetamiento}\label{tamauxf1o-del-empaquetamiento}

El empaquetamiento combinado tiene tamaño

\[
\begin{aligned}
\nu_3(G)
&\ge
|F|+\frac{\binom p2-|F|-O(p)}3\\
&=
\frac13\binom p2+\frac23|F|-O(p)\\
&\ge
\frac13\binom p2+\frac13\sum_i d_i-O(p+q).
\end{aligned}
\tag{8.10}
\]

En consecuencia,

\[
\begin{aligned}
\Phi(G)
&\le
\frac13\binom p2+\frac13\sum_i d_i+O(p+q)\\
&\le
\frac13\binom p2+\frac{pq}{3}+O(p+q)\\
&=
\frac{(p+q)^2}{6}-\frac{p+q^2}{6}+O(p+q).
\end{aligned}
\tag{8.11}
\]

Esto es a lo sumo \(n^2/6+O(n)\).

\begin{center}\rule{0.5\linewidth}{0.5pt}\end{center}

\section{9. Demostración del teorema
principal}\label{demostraciuxf3n-del-teorema-principal}

Demostramos ahora el Teorema 1.1 por contradicción.

Supongamos que ninguna constante absoluta funciona. Para cada entero
positivo \(k\), elegimos un grafo split \(G_k\), mínimo en su número
\(n_k\) de vértices, tal que

\[
\Phi(G_k)
>
\frac{n_k^2}{6}+kn_k.
\tag{9.1}
\]

Necesariamente \(n_k\to\infty\). Para todo \(v\in I(G_k)\), la
minimalidad da

\[
\Phi(G_k)
\le
\Phi(G_k-v)+d(v),
\]

y por tanto

\[
\boxed{
d(v)>
\frac{2n_k-1}{6}+k.
}
\tag{9.2}
\]

Primero, si \(q_k\ge2p_k-1\), el Lema 5.1 da

\[
\Phi(G_k)
\le
\frac{n_k^2}{6}+\frac{p_k}{2},
\]

lo que contradice (9.1) para \(k\) grande. Podemos entonces suponer
\(q_k<2p_k-1\) y pasar a una subsucesión según

\[
\alpha_k=\frac{q_k}{p_k}\in[0,2).
\]

\subsection{9.1 Régimen bulk}\label{ruxe9gimen-bulk}

Supongamos que para algún \(\varepsilon>0\),

\[
\varepsilon\le\alpha_k\le2-\varepsilon.
\]

El Teorema 4.2 da

\[
\nu_3^*(G_k)
\ge
T(G_k)+c_\varepsilon p_k^2-O(p_k).
\]

Haxell--Rödl da

\[
\nu_3(G_k)
\ge
\nu_3^*(G_k)-o(n_k^2).
\]

Para \(k\) suficientemente grande, el margen cuadrático domina la
pérdida de integralidad, de modo que \(\nu_3(G_k)\ge T(G_k)\), en
contradicción con (9.1).

\subsection{\texorpdfstring{9.2 El extremo
\(\alpha_k\to0\)}{9.2 El extremo \textbackslash alpha\_k\textbackslash to0}}\label{el-extremo-alpha_kto0}

Este caso se cierra mediante la Sección 8, que da una constante absoluta
\(C_0\) tal que

\[
\Phi(G_k)
\le
\frac{n_k^2}{6}+C_0n_k.
\]

Para \(k>C_0\), esto contradice (9.1).

\subsection{\texorpdfstring{9.3 El extremo
\(\alpha_k\to2\)}{9.3 El extremo \textbackslash alpha\_k\textbackslash to2}}\label{el-extremo-alpha_kto2}

Escribimos

\[
q=2p-s,
\qquad
s=o(p).
\]

Si \(s=O(\sqrt p)\), la desigualdad (5.3) da la cota de error lineal
requerida. Supongamos desde ahora

\[
\sqrt p\ll s=o(p).
\tag{9.3}
\]

Para todos los miembros suficientemente grandes de la sucesión,

\[
p\ge2304,
\qquad
6\sqrt p\le s\le\frac p8,
\qquad
k\ge1.
\tag{9.4}
\]

Por (9.2),

\[
m_i=|S_i|
<
\frac s3+\frac16-k
\le
\frac s3-\frac56.
\]

Como cada \(m_i\) es entero, si \(m=\max_i m_i\), entonces

\[
\boxed{3m\le s-3.}
\tag{9.5}
\]

Para \(x\in K\), sea

\[
a_x=|\{i:x\in S_i\}|
\]

y

\[
D=\sum_{x\in K}a_x(q-a_x).
\tag{9.6}
\]

Entonces

\[
\sum_{i,j}|S_i\triangle S_j|=2D.
\tag{9.7}
\]

\subsubsection{Dispersión alta}\label{dispersiuxf3n-alta}

Supongamos

\[
D\ge\frac{qs^2}{12}.
\tag{9.8}
\]

Por (9.5),

\[
2p-3m-1\ge2p-s+2=q+2.
\]

El Lema 6.1 y (9.7) dan entonces

\[
V
\ge
\frac{q+2}{2}D
\ge
\frac q2D.
\tag{9.9}
\]

El coeficiente de doble factor del Lema 5.2 satisface

\[
\boxed{\delta\ge\frac78.}
\tag{9.10}
\]

En efecto, si \(p\) es impar, entonces \(r_p=p\) y

\[
\delta=\frac{p-s}{p}\ge\frac78,
\]

mientras que si \(p\) es par, entonces \(r_p=p-1\) y

\[
\delta=\frac{p+1-s}{p-1}\ge\frac78.
\]

Además, \(S_2\ge M^2/q\), y (9.5) implica \(M/q\le m\le(s-3)/3\). Como
la parábola resultante es creciente en este intervalo,

\[
\frac{(s-1)M-S_2}{q}
\le
\frac{2s(s-3)}9.
\tag{9.11}
\]

Al sustituir (9.8)--(9.11) en el Lema 5.2, obtenemos

\[
\Phi(G)-\frac{n^2}{6}
\le
\frac p2-\frac{5s^2}{288}-\frac{2s}{3}.
\tag{9.12}
\]

Como \(s^2\ge36p\),

\[
\Phi(G)-\frac{n^2}{6}
\le
-\frac p8-\frac{2s}{3}<0,
\]

una contradicción.

\subsubsection{Dispersión baja}\label{dispersiuxf3n-baja}

Supongamos

\[
D<\frac{qs^2}{12}.
\tag{9.13}
\]

Por (9.7), existe algún \(j\) tal que

\[
\sum_i|S_i\triangle S_j|
<
\frac{s^2}{6}.
\tag{9.14}
\]

Definimos

\[
R=S_j,
\qquad
\rho=|R|.
\]

Entonces

\[
A_R+B_R<\frac{s^2}{6},
\qquad
\rho\le m\le\frac{s-3}{3}.
\tag{9.15}
\]

Verificamos explícitamente las hipótesis del Lema 7.1. Como
\(s\le p/8\),

\[
q=2p-s\ge\frac{15p}{8},
\qquad
b=p-\rho\ge p-\frac s3.
\]

Por tanto \(b\ge2\), \(q\ge r_b\) y \(b\ge\chi'(K_\rho)\). Además
\(r_b\ge b-1\), luego

\[
u=q-r_b\le p-s+\rho+1.
\]

Para todo \(i\),

\[
2\rho+t_i+1
\le
3m+1
\le
s-2.
\]

En consecuencia,

\[
b-t_i\ge\max\{\rho,u\},
\]

y el Lema 7.1 se aplica.

Si \(\rho=0\), entonces \(B_R=0\). Si \(\rho\ge1\), entonces

\[
\theta_R
\le
\frac{\rho}{p-\rho}
\le
\frac{8s}{23p}.
\]

Usando \(u=q-r_b\), \(r_b\le b\) y \(q\ge15p/8\), obtenemos

\[
\kappa_R
=1-2(1-\theta_R)\frac uq
\le
\frac sq+2\theta_R
\le
\frac{5s}{4p}.
\tag{9.16}
\]

Del mismo modo,

\[
\frac{(s-2\rho-1)_+}{q}
\le
\frac{5s}{4p}.
\tag{9.17}
\]

Por (9.15), la desviación positiva total de (7.6) es entonces a lo sumo

\[
\frac{5s}{4p}(A_R+B_R)
<
\frac{5s^3}{24p}
\le
\frac{5s^2}{192}.
\tag{9.18}
\]

Mientras tanto,

\[
\frac{s^2}{6}-s\rho+2\rho^2
=
2\left(\rho-\frac s4\right)^2+\frac{s^2}{24}
\ge
\frac{s^2}{24}.
\tag{9.19}
\]

El Lema 7.1 da ahora

\[
\Phi(G)-\frac{n^2}{6}
\le
\frac p2-\frac{s^2}{64}.
\tag{9.20}
\]

Como \(s^2\ge36p\), el lado derecho es a lo sumo \(-p/16\), nuevamente
una contradicción.

Todas las subsucesiones posibles son imposibles. Se sigue el Teorema
1.1. \(\square\)

\begin{center}\rule{0.5\linewidth}{0.5pt}\end{center}

\section{10. Corolarios}\label{corolarios}

\subsection{10.1 Particiones en cliques}\label{particiones-en-cliques}

El Corolario 1.2 se sigue inmediatamente de

\[
\operatorname{cp}(G)
\le
|E(G)|-2\nu_3(G).
\]

\subsection{10.2 Optimalidad del término
principal}\label{optimalidad-del-tuxe9rmino-principal}

La familia

\[
K_p\vee\overline K_{2p}
\]

tiene \(n=3p\). Una factorización de \(K_p\) asignada a vértices
independientes empaqueta toda arista de la clique en un triángulo
\(KKI\). La expresión resultante del empaquetamiento es

\[
|E|-2\nu_3
=
\frac{n^2}{6}+\frac n6.
\tag{10.1}
\]

La construcción de Erdős--Ordman--Zalcstein para particiones en cliques
sobre esta familia muestra que el coeficiente \(1/6\) no puede
mejorarse.

\subsection{10.3 Localización de sucesiones
difíciles}\label{localizaciuxf3n-de-sucesiones-difuxedciles}

La Proposición 1.3 registra la localización proporcionada por la
demostración. Debe entenderse como una descripción estructural de un
fallo hipotético de toda estimación con error lineal, no como un teorema
extremal adicional posterior al Teorema 1.1. El margen bulk excluye
razones acotadas lejos de \(0\) y \(2\), la construcción escasa excluye
\(|I|/|K|\to0\), y la factorización promediada excluye el corredor corto
cerca de \(|I|=2|K|\). La única ubicación restante es el corredor
mesoscópico

\[
|I|=2|K|-s,
\qquad
\sqrt{|K|}\ll s=o(|K|),
\]

que se elimina mediante la dicotomía de dispersión.

\subsection{10.4 Valor fraccionario exacto de perfil
común}\label{valor-fraccionario-exacto-de-perfil-comuxfan}

El Teorema 3.1 es útil por sí mismo: da el empaquetamiento fraccionario
exacto de triángulos para todo grafo split en el que todos los vértices
independientes tienen un vecindario común, incluido el régimen con
demasiado pocos vértices independientes para colorear integralmente
todas las aristas del vecindario. Lo enunciamos como corolario
independiente.

\subsubsection{Corolario 10.4}\label{corolario-10.4-empaquetamiento-exacto-con-vecindario-comuxfan}

Para todos los enteros \(p\ge3\), \(q\ge0\) y \(0\le d\le p\), el grafo
\(H(p,q,d)\) ---clique de orden \(p\); \(q\) vértices independientes,
cada uno con el mismo vecindario de tamaño \(d\)--- satisface

\[
\nu_3^*(H(p,q,d))=F(p,q,d)=\min\left\{\tfrac{\binom p2+qd}{3},\ \binom d2+\binom{p-d}2,\ \binom d2+\tfrac{d(p-d)+\binom{p-d}2}{3}\right\}.
\]

Esta es una fórmula cerrada exacta para el número fraccionario de
empaquetamiento de triángulos de una familia completa de un parámetro,
válida para todo \(d\), incluido el régimen de \(q\) pequeño donde
ninguna factorización integral cubre el vecindario; puede tener interés
independiente fuera del problema de particiones en cliques.

\subsubsection{Corolario 10.4b}\label{corolario-10.4b-grafos-threshold}

Como \(K_p\vee\overline K_{2p}\) es un grafo threshold y todo grafo
threshold es split, el Teorema 1.1 da, en particular, la cota de error
lineal

\[
\operatorname{cp}(G)\le n^2/6+Cn
\]

para todos los grafos threshold, con la misma constante principal óptima
atestiguada dentro de esa subclase.

\subsection{\texorpdfstring{10.5 Efectividad localizada cerca de
\(q=2p\)}{10.5 Efectividad localizada cerca de q=2p}}\label{efectividad-localizada-cerca-de-q2p}

\subsubsection{Proposición 10.1}\label{proposiciuxf3n-10.1-cotas-expluxedcitas-para-el-corredor}

Sea \(q=2p-s\) y \(n=p+q\).

\begin{enumerate}
\def\labelenumi{\arabic{enumi}.}
\item
  Si \[
  p\ge36,
  \qquad
  0\le s\le6\sqrt p,
  \] entonces \[
  \boxed{
  \Phi(G)\le\frac{n^2}{6}+2n.
  }
  \tag{10.2}
  \]
\item
  Si \[
  p\ge2304,
  \qquad
  6\sqrt p\le s\le\frac p8,
  \] y todo \(v\in I\) satisface \[
  d(v)>\frac{2n-1}{6}+1,
  \] entonces \[
  \boxed{
  \Phi(G)\le\frac{n^2}{6}.
  }
  \tag{10.3}
  \]
\end{enumerate}

En consecuencia, para todo \(C\ge2\), un contraejemplo de orden mínimo a

\[
\Phi(G)\le\frac{n^2}{6}+Cn
\]

no puede satisfacer

\[
p\ge2304,
\qquad
0\le s\le\frac p8.
\]

\subsubsection{Demostración}\label{demostraciuxf3n-6}

Para el corredor corto, (5.3) y \(s^2\le36p\) dan

\[
\Phi(G)-\frac{n^2}{6}
\le
\frac{7p}{2}+\frac14
\le
2n,
\]

porque \(p\ge36\) implica \(s\le p\) y, por tanto, \(n=3p-s\ge2p\).

La afirmación mesoscópica es exactamente el cálculo cuantitativo de
dispersión alta/baja de la Sección 9.3. La consecuencia final se sigue
de la desigualdad de grado obtenida al eliminar un vértice independiente
de un contraejemplo de orden mínimo. \(\square\)

\begin{center}\rule{0.5\linewidth}{0.5pt}\end{center}

\section{11. Discusión}\label{discusiuxf3n}

\subsection{11.1 De la asintótica fraccionaria a un teorema integral con
error
lineal}\label{de-la-asintuxf3tica-fraccionaria-a-un-teorema-integral-con-error-lineal}

La dificultad central es que el teorema general

\[
\nu_3^*(G)-\nu_3(G)=o(n^2)
\]

no proporciona una tasa lineal. La presente demostración evita exigir un
único mecanismo uniforme de redondeo.

\begin{itemize}
\tightlist
\item
  En el bulk, el LP de perfil común produce holgura cuadrática.
\item
  Cerca de \(q=2p\), las estructuras explícitas de factorización
  reemplazan el redondeo general.
\item
  Cerca de \(q=0\), la clique residual se vuelve suficientemente densa
  para una descomposición exacta.
\end{itemize}

Esta estrategia dependiente del régimen es más débil que un teorema
universal de brecha lineal de integralidad, pero basta para el problema
extremal de particiones en cliques.

\subsection{11.2 Por qué el grado máximo residual no es el invariante
correcto}\label{por-quuxe9-el-grado-muxe1ximo-residual-no-es-el-invariante-correcto}

Los primeros trabajos exploratorios intentaron probar que un
empaquetamiento óptimo deja un residuo de clique de grado máximo
acotado. El comportamiento numérico dependía fuertemente del algoritmo
de empaquetamiento, y la demostración final no necesita ninguna
propiedad de este tipo.

Los invariantes efectivos son, en cambio:

\begin{itemize}
\tightlist
\item
  el margen fraccionario de perfil común;
\item
  el término de varianza \(S_2\);
\item
  la energía de polarización \(V\);
\item
  las desviaciones centradas \(A_R,B_R\);
\item
  la divisibilidad del residuo denso.
\end{itemize}

\subsection{11.3 Efectividad localizada y no efectividad
restante}\label{efectividad-localizada-y-no-efectividad-restante}

La Proposición 10.1 muestra que todo el corredor casi extremal tratado
mediante factorización, polarización y completación respecto de un
centro desplazado es cuantitativamente efectivo. Se puede tomar

\[
C_{\mathrm{corr}}=2,
\qquad
p_0=2304,
\qquad
s_0(p)=6\sqrt p,
\qquad
\eta_0=\frac18.
\]

La constante global \(C\) sigue siendo no efectiva por dos razones
independientes.

\begin{enumerate}
\def\labelenumi{\arabic{enumi}.}
\tightlist
\item
  El teorema de Haxell--Rödl se usa mediante una afirmación asintótica
  \(o(n^2)\) en el bulk.
\item
  Los teoremas de descomposición densa en triángulos se invocan con
  umbrales no especificados de «suficientemente grande» en el régimen
  escaso.
\end{enumerate}

Así, no queda no efectividad oculta en las Secciones 5--7; proviene
únicamente de los dos insumos asintóticos externos anteriores. Además,
después del Apéndice D, las Secciones 5--7 son autocontenidas: la
Proposición 10.1 no depende de ningún teorema externo.

\subsection{11.4 Relación con el teorema extremal
fraccionario}\label{relaciuxf3n-con-el-teorema-extremal-fraccionario}

El teorema extremal complete-split determina el máximo fraccionario
global de

\[
|E(G)|-2\nu_3^*(G)
\]

tanto sobre grafos split como sobre grafos cordales. El LP de perfil
común del presente paper cumple otro propósito: es una cota inferior
local, sensible al perfil, para el grafo replicado, y su holgura
cuantitativa absorbe la pérdida de integralidad fuera del corredor
extremal.

Por tanto, el cálculo extremal fraccionario global y el cálculo local de
perfil común son compatibles, pero lógicamente distintos. La presente
demostración es autocontenida salvo por los teoremas externos enunciados
en la Sección 2.

\textbf{Observación (clonado a través de la serie).} El Lema 4.1 y el
paso de simetrización por clones del Paper II usan la misma operación
subyacente en funciones distintas. El Paper II usa clonado por pares
para llevar un grafo cordal hacia un grafo terminal complete-split. El
Lema 4.1 usa clonado \(q\)-fold en el lado de cobertura para comparar un
perfil arbitrario con un grafo de perfil común. Se conserva entonces el
término compartido \emph{clonado}, mientras que el calificativo
\emph{fraccional} distingue la desigualdad presente.

\subsection{11.5 Relación con el problema cordal de error
lineal}\label{relaciuxf3n-con-el-problema-cordal-de-error-lineal}

La reducción complete-split establece la cota cordal asintóticamente
óptima

\[
\operatorname{cp}(G)
\le
\left(\frac16+o(1)\right)n^2.
\]

La estimación más fuerte

\[
\operatorname{cp}(G)
\le
\frac{n^2}{6}+O(n)
\]

permanece abierta para grafos cordales generales. Los mecanismos
desarrollados aquí sugieren un programa concreto de tres partes.

\begin{enumerate}
\def\labelenumi{\arabic{enumi}.}
\item
  \textbf{Estabilidad del teorema extremal fraccionario.} Reforzar el
  teorema fraccionario cordal exacto del Paper II hasta una afirmación
  de estabilidad: todo grafo cordal cuyo valor está a distancia
  \(\delta n^2\) del máximo fraccionario está a distancia de edición
  \(\varepsilon n^2\) de un grafo complete-split maximizador. Fuera de
  la casi extremalidad, la holgura fraccionaria cuadrática absorbe la
  pérdida general de integralidad exactamente como en la Sección 4.
\item
  \textbf{Mecanismos robustos del corredor.} Extender los argumentos de
  factorización, polarización y centro desplazado de las Secciones 5--7
  desde grafos split a perturbaciones de \(\varepsilon n^2\) de grafos
  complete-split, para aplicarlos a los grafos casi extremales
  producidos por estabilidad.
\item
  \textbf{Ensamblaje mediante el árbol de cliques.} Parchar las
  estimaciones locales a través del árbol de cliques mediante un ledger
  de propiedad de aristas que preserve las capacidades de los
  separadores, en el espíritu de la completación respecto de un centro
  desplazado, de modo que ninguna arista se cargue dos veces.
\end{enumerate}

Una demostración cordal con error lineal todavía tendría que preservar
la propiedad y las capacidades de los separadores a lo largo del árbol
de cliques; el ledger del paso 3 está pensado precisamente para ese
propósito.

Un paper compañero en preparación investiga este programa. Su recursión
prevista por separadores-clique usa el teorema split demostrado aquí
como caso base y está diseñada para que los regímenes cordales reflejen
los ingredientes correspondientes de las Secciones 3--9: evaluación de
perfil común, clonado fraccional, factorización y análisis del corredor.
Hasta que ese argumento compañero esté completo, este párrafo registra
la interfaz propuesta y no un teorema del presente paper.

\subsection{11.6 Perímetro y estado actual de la verificación
formal}\label{peruxedmetro-y-estado-actual-de-la-verificaciuxf3n-formal}

Los Papers I y II de la serie están formalmente verificados en Lean 4
con Mathlib. El Paper III se está verificando bajo el mismo protocolo de
liberación. El proyecto Lean del presente manuscrito usa Mathlib v4.28.0
y compila actualmente como un grafo completo de dependencias.

El desarrollo formal separa dos capas.

\begin{itemize}
\tightlist
\item
  \textbf{Capa X} contiene exactamente los dos insumos asintóticos
  externos, los Teoremas 2.1 y 2.3, retenidos como axiomas nombrados con
  sus hipótesis enunciadas explícitamente.
\item
  \textbf{Capa E} contiene el núcleo matemático finito: Teoremas 3.1 y
  4.2; Lemas 4.1--7.1; la corrección de divisibilidad del Apéndice B; el
  desarrollo de coloreo de aristas por listas del Apéndice D; las
  cadenas cuantitativas de las Secciones 9.3 y 10.5; la Proposición
  10.1; y el ensamblaje final.
\end{itemize}

A la fecha de este borrador, el proyecto tiene ocho obligaciones
localizadas de tipo \texttt{sorry}. Están confinadas a: los pequeños
casos degenerados del LP de perfil común; el bloque de redondeo con uno
y dos factores de la Sección 5; la conclusión del corredor medio de la
Proposición 10.1; un helper de matching en el desarrollo autocontenido
de Galvin/König; la construcción de empaquetamiento con tres familias
que subyace a la Sección 7.2; y la estimación de empaquetamiento del
régimen escaso que invoca el Teorema 2.3. El enunciado y el ensamblaje
de dependencias del Teorema 1.1 ya compilan, pero el teorema permanece
transitivamente abierto hasta cerrar estas hojas.

Ningún cálculo ni script de regresión es una premisa de la demostración.
Las comprobaciones de aritmética exacta y programación entera son
auditorías suplementarias. Bajo una convención única de conteo,
comprenden 46,390 comprobaciones exactas de aritmética y LP, y 91
comprobaciones exactas de ILP, para un total de 46,481, sin
discrepancias reportadas.

La línea de estado del preprint público se actualizará solo después de
congelar todos los siguientes elementos: un build con cero
\texttt{sorry}, el informe de axiomas, el toolchain y manifest de Lean,
y el commit de liberación. Por tanto, la afirmación final prevista ya
está delimitada, pero no se declara en este borrador.

\begin{center}\rule{0.5\linewidth}{0.5pt}\end{center}

\section{12. Usos potenciales y direcciones
futuras}\label{usos-potenciales-y-direcciones-futuras}

\subsection{12.1 Brecha universal lineal de
integralidad}\label{brecha-universal-lineal-de-integralidad}

\subsubsection{Problema 12.1}\label{problema-12.1}

¿Existe una constante absoluta \(C\) tal que

\[
\nu_3^*(G)-\nu_3(G)
\le C|V(G)|
\]

para todo grafo split \(G\)?

El presente paper proporciona estructuras de apoyo cerca del corredor
extremal, pero no resuelve el régimen bulk sin holgura cuadrática.

\subsection{12.2 Constantes globales
efectivas}\label{constantes-globales-efectivas}

\subsubsection{Problema 12.2}\label{problema-12.2}

Encontrar una constante global explícita \(C\) en el Teorema 1.1.

El corredor casi extremal ya es efectivo por la Proposición 10.1. Un
teorema global completamente cuantitativo requiere entonces solo tasas
explícitas en el paso fraccionario-a-integral, o un reemplazo
estructurado de Haxell--Rödl en el bulk, junto con umbrales efectivos de
descomposición densa en el régimen con lado independiente escaso.

\subsection{12.3 Empaquetamiento
algorítmico}\label{empaquetamiento-algoruxedtmico}

La demostración contiene piezas implementables en tiempo polinomial:

\begin{itemize}
\tightlist
\item
  factorizaciones de grafos completos;
\item
  selección ponderada de vértices reservados;
\item
  coloreo bipartito de aristas por listas;
\item
  extracción de matchings;
\item
  algoritmos de descomposición densa implícitos en la absorción
  iterativa.
\end{itemize}

Es natural buscar un algoritmo polinomial que produzca una partición en
cliques de tamaño

\[
\frac{n^2}{6}+O(n)
\]

con una constante explícita. Los pasos de factorización, coloreo por
listas y matchings de las Secciones 5--7 son individualmente
polinomiales; ensamblarlos en un solo algoritmo de aproximación con
garantía demostrada queda para trabajo futuro.

\subsubsection{Corolario 12.2}\label{corolario-12.2-algoritmo-efectivo-en-el-corredor}

En el corredor casi extremal cubierto por la Proposición 10.1 ---con
\(p_0,s_0\) explícitos---, la demostración es constructiva: las
factorizaciones simples y dobles de \(K_p\), la selección de vértices
reservados, el coloreo de aristas por listas ---Apéndice D, cuya
demostración mediante núcleos y emparejamientos estables es a su vez un
algoritmo--- y la extracción de matchings se ejecutan todos en tiempo
polinomial. Por tanto, existe un algoritmo determinista de tiempo
polinomial que, para grafos split en el corredor efectivo, produce una
partición en cliques de tamaño \(n^2/6+2n\), con la constante explícita
de la Proposición 10.1. Fuera del corredor, la garantía actual depende
del insumo no constructivo del bulk ---Teorema 2.1---; un algoritmo
global completamente efectivo queda para trabajo futuro.

\subsection{12.4 Estabilidad y clasificación
extremal}\label{estabilidad-y-clasificaciuxf3n-extremal}

\subsubsection{Problema 12.3}\label{problema-12.3}

Caracterizar los grafos split que satisfacen

\[
\operatorname{cp}(G)
\ge
\frac{n^2}{6}-o(n^2).
\]

La demostración sugiere que los grafos casi extremales deben satisfacer
\(|I|=(2+o(1))|K|\) y tener perfiles de ausencias de baja dispersión
cercanos a un centro común.

Dos evidencias apoyan una afirmación de estabilidad de tipo Simonovits.
Primero, una enumeración exhaustiva para \(n=9\) ---todos los grafos
split con \(|K|=3\)--- muestra que los grafos a distancia aditiva
constante del máximo están a distancia de edición acotada de la familia
complete-split, mientras que los perfiles genuinamente dispersos quedan
una cantidad lineal por debajo del máximo. Segundo, el margen
fraccionario del Teorema 4.2 ya da una cota de conjunto de nivel: lejos
de la casi extremalidad, la holgura cuadrática \(\mu(\alpha)p^2\) es
estrictamente positiva, por lo que todo casi maximizador debe tener
\(\alpha\) cerca de la razón extremal y pequeña dispersión de perfiles.
Un teorema completo de estabilidad aún requeriría cuantificar la
ganancia de un único paso de simetrización y controlar las direcciones a
lo largo de las cuales \(\Phi^*\) es plana; lo dejamos como problema
abierto, formulado de manera natural en la métrica de corte
---graphon--- para absorber la elección del grafo complete-split más
cercano.

\subsection{12.5 Empaquetamiento con cliques mayores y error
lineal}\label{empaquetamiento-con-cliques-mayores-y-error-lineal}

La reducción complete-split determina el problema extremal fraccionario
y la asintótica integral de primer orden para particiones en \(K_r\)
fijos y aristas individuales. El análogo natural para cliques mayores es
la siguiente pregunta de error lineal.

Para \(r\ge3\) fijo, ¿satisface todo grafo split

\[
\pi_r(G)
\le
\frac{r-1}{4r}n^2+O_r(n),
\]

donde \(\pi_r(G)\) es el número mínimo de partes en una partición de
aristas en copias de \(K_r\) y aristas individuales?

Las ideas de perfil común y clonado sugieren una posible vía, pero los
programas locales por órbitas y los mecanismos de redondeo integral
relevantes se vuelven sustancialmente más complicados. Los métodos de
diseños y descomposición, como los desarrollados por Keevash {[}10{]} y
el marco de rebanadas regulares de Allen, Böttcher, Cooley y Mycroft
{[}11{]}, pueden ser herramientas útiles; no son ingredientes de la
presente demostración.

\subsection{12.6 Aplicaciones cordales y a árboles de
cliques}\label{aplicaciones-cordales-y-a-uxe1rboles-de-cliques}

La completación respecto de un centro desplazado fue diseñada para
separar aristas propias de la clique de aristas protegidas de interfaz.
Una versión paramétrica puede ser útil en un futuro teorema cordal con
error lineal, donde los separadores son compartidos entre cliques
maximales vecinas.

\subsection{12.7 Principio reutilizable de
demostración}\label{principio-reutilizable-de-demostraciuxf3n}

El paper ilustra una estrategia más amplia:

\begin{quote}
Usar LP simétricos exactos para crear holgura lejos de la extremalidad,
y reservar el redondeo combinatorio explícito solo para el corredor
donde la holgura desaparece.
\end{quote}

Este principio puede aplicarse a otros problemas estructurados de
empaquetamiento y cobertura en los que un teorema general de regularidad
es demasiado débil en segundo orden.

\begin{center}\rule{0.5\linewidth}{0.5pt}\end{center}

\section{13. Reproducibilidad}\label{reproducibilidad}

El paquete de liberación de trabajo para este borrador se organiza como
sigue.

\begin{Shaded}
\begin{Highlighting}[]
\NormalTok{01\_MANUSCRIPT/}
\NormalTok{    PAPER\_III\_split\_lineal\_v0.9.4\_review\_es.md}
\NormalTok{02\_PROOF\_CONTROL/}
\NormalTok{    LEDGER.md}
\NormalTok{    LEDGER\_INCREMENTAL\_v0.9.2\_formalization\_draft.md}
\NormalTok{    AUDIT\_PAPER\_C.md}
\NormalTok{03\_FORMALIZATION/}
\NormalTok{    lean/                         (fuentes Lean 4 / Mathlib v4.28.0)}
\NormalTok{    lakefile.toml}
\NormalTok{    lean{-}toolchain}
\NormalTok{    lake{-}manifest.json}
\NormalTok{    BUILD\_STATUS.md               (temporal hasta congelar el build de liberación)}
\NormalTok{04\_SUPPLEMENTARY\_REGRESSION/}
\NormalTok{    audit\_c\_fast.py}
\NormalTok{    audit\_c\_ilp.py}
\NormalTok{    audit\_c\_ilp\_results.txt}
\end{Highlighting}
\end{Shaded}

El manuscrito es la fuente matemática de verdad. \texttt{LEDGER.md} es
la especificación autoritativa de dependencias para el desarrollo Lean;
el ledger incremental registra únicamente la transición desde v0.9.1 a
este borrador. El paquete final de liberación añadirá el commit
congelado, el log de build, el informe de axiomas y el manifiesto
SHA-256.

La demostración misma está contenida en el manuscrito. La optimización
de perfil común se prueba en el Teorema 3.1; el margen fraccionario se
prueba en el Teorema 4.2 y el Apéndice A; las desigualdades de
factorización, polarización y centro desplazado se prueban en los Lemas
5.1--7.1; la corrección de divisibilidad se prueba en la Sección 8.3 y
el Apéndice B; y las constantes explícitas casi extremales se prueban
algebraicamente en las Secciones 9.3 y 10.5.

El script incluido es una prueba suplementaria de regresión para las
constantes explícitas del corredor. Su salida no se usa como hipótesis,
reducción ni paso de demostración en ninguna parte del manuscrito. Los
cálculos exploratorios históricos no forman parte del paquete lógico de
la demostración.

\begin{center}\rule{0.5\linewidth}{0.5pt}\end{center}

\section{Apéndice A. Álgebra del margen
fraccionario}\label{apuxe9ndice-a.-uxe1lgebra-del-margen-fraccionario}

Sea \(x=d/p\) y \(\alpha=q/p\). Ignorando temporalmente los términos
lineales exactos, las tres ramas normalizadas de \(F\) son

\[
f_1(x)=\frac16+\frac{\alpha x}{3},
\]

\[
f_2(x)=x^2-x+\frac12,
\]

y

\[
f_3(x)=\frac16+\frac{x^2}{3}.
\]

Al restar la contribución afín objetivo se obtienen los mínimos

\[
\min_x\left(f_1(x)-\frac{\alpha x}{2}-C_\alpha\right)
=
\frac{\alpha^2}{12},
\]

\[
\min_x\left(f_2(x)-\frac{\alpha x}{2}-C_\alpha\right)
=
\frac{(2-\alpha)^2}{48},
\]

y

\[
\min_x\left(f_3(x)-\frac{\alpha x}{2}-C_\alpha\right)
=
\begin{cases}
\alpha(8-5\alpha)/48,&\alpha\le4/3,\\
(2-\alpha)^2/12,&\alpha\ge4/3.
\end{cases}
\]

La tercera rama está dominada por la envolvente inferior de las dos
primeras. Al restituir la contribución exacta \(-p/2\) en cada
estimación puntual y promediar, se produce el término \(-p/4\) del
Teorema 4.2.

\begin{center}\rule{0.5\linewidth}{0.5pt}\end{center}

\section{Apéndice B. Corrección de
divisibilidad}\label{apuxe9ndice-b.-correcciuxf3n-de-divisibilidad}

Sea \(P=x_1\cdots x_p\) un camino y sea \(O\subseteq V(P)\) de
cardinalidad par. Definimos

\[
J=\{x_jx_{j+1}:|O\cap\{x_1,\ldots,x_j\}|\text{ es impar}\}.
\]

Todo vértice interno \(x_j\) cambia la paridad del prefijo exactamente
cuando \(x_j\in O\). Por tanto,

\[
\operatorname{Odd}(J)=O.
\]

Además, \(|E(J)|\le p-1\) y \(\Delta(J)\le2\).

Después de corregir la paridad, el grado mínimo ha disminuido a lo sumo
en dos. En el régimen con lado independiente escaso sigue siendo mayor
que \(3p/4\) para todo \(p\) suficientemente grande, de modo que el
teorema de Turán proporciona un \(K_5\). Eliminar un \(C_4\) dentro de
ese \(K_5\) cambia el número de aristas en \(1\pmod3\), mientras que
eliminar un \(C_5\) lo cambia en \(2\pmod3\). Todo grado afectado cambia
en dos, por lo que la paridad se preserva. La pérdida total de grado
debida al subgrafo del camino y al ciclo corrector es a lo sumo cuatro.
Así, cuando \(q=o(p)\), el grafo final tiene grado mínimo al menos
\(p-1-q-4\), que eventualmente excede \((0.9+\varepsilon_0)p\) para todo
\(0<\varepsilon_0<0.1\) fijo.

\begin{center}\rule{0.5\linewidth}{0.5pt}\end{center}

\section{Apéndice C. Auditorías
computacionales}\label{apuxe9ndice-c.-auditoruxedas-computacionales}

\subsection{C.1 LP exacto de perfil
común}\label{c.1-lp-exacto-de-perfil-comuxfan}

\texttt{verify\_common\_profile\_lp.py} enumera los vértices del
poliedro simetrizado de coberturas mediante eliminación gaussiana
racional exacta y compara el óptimo con (3.5).

\subsection{C.2 Margen fraccionario
exacto}\label{c.2-margen-fraccionario-exacto}

\texttt{verify\_fractional\_margin.py} comprueba el Teorema 4.2 para

\[
3\le p\le80,
\qquad
1\le q\le2p,
\qquad
0\le d\le p
\]

usando aritmética racional exacta.

\subsection{C.3 Auditorías ILP
pequeñas}\label{c.3-auditoruxedas-ilp-pequeuxf1as}

\texttt{verify\_factor\_rounding.py} y
\texttt{verify\_shifted\_center.py} calculan empaquetamientos integrales
exactos de triángulos en instancias pequeñas mediante un ILP binario de
capacidad de aristas. Verifican las cotas superiores declaradas para
\(\Phi(G)\).

\subsection{C.4 Polarización}\label{c.4-polarizaciuxf3n}

\texttt{verify\_polarization.py} realiza comprobaciones exhaustivas de
perfiles pequeños y comprobaciones aleatorias con enteros exactos del
Lema 6.1.

\subsection{C.5 Divisibilidad}\label{c.5-divisibilidad}

\texttt{verify\_divisibility\_correction.py} verifica exhaustivamente la
construcción de paridad sobre caminos hasta orden dieciocho, comprueba
exactamente el umbral fijo de grado mínimo y sigue la corrección
completa de paridad y módulo tres sobre residuos densos aleatorios.

\begin{center}\rule{0.5\linewidth}{0.5pt}\end{center}

\section{Apéndice D. Coloreo de aristas por listas
autocontenido}\label{apuxe9ndice-d.-coloreo-de-aristas-por-listas-autocontenido}

Este apéndice demuestra el teorema de coloreo de aristas por listas
usado en la Sección 7.2 ---enunciado como Teorema 2.2---, de modo que el
argumento del corredor no dependa de ningún teorema externo de coloreo.
El resultado es el caso de grado máximo del teorema de Galvin {[}12{]};
el grafo de ganancias de la Sección 7.2 es simple y bipartito, que es el
caso probado aquí.

\subsection{D.1 Núcleos}\label{d.1-nuxfacleos}

Un \textbf{núcleo} de un dígrafo \(D\) es un conjunto independiente
\(K\) de vértices tal que todo vértice fuera de \(K\) tiene un vecino de
salida en \(K\). El dígrafo \(D\) es \textbf{perfecto por núcleos} si
todo subdígrafo inducido tiene un núcleo. Un núcleo de un dígrafo no
vacío es no vacío.

\subsubsection{Lema D.1}\label{lema-d.1-lema-de-coloreo-por-nuxfacleos}

Sea \(D\) un dígrafo perfecto por núcleos en el que cada vértice \(v\)
tiene una lista \(L(v)\) con \(|L(v)|\ge d^+_D(v)+1\). Entonces el grafo
subyacente de \(D\) admite un coloreo propio que elige cada color desde
la lista correspondiente.

\subsubsection{Demostración}\label{demostraciuxf3n-7}

Elegimos cualquier color \(c\) que aparezca en alguna lista, definimos
\(S=\{v:c\in L(v)\}\) y sea \(K\) un núcleo de \(D[S]\). Coloreamos
todos los vértices de \(K\) con \(c\); el coloreo es propio sobre \(K\)
porque \(K\) es independiente. Eliminamos \(K\) de \(D\) y eliminamos
\(c\) de la lista de todo vértice de \(S\setminus K\). Cada
\(v\in S\setminus K\) perdió un color, pero también al menos un vecino
de salida, a saber, su vecino de salida en \(K\), de modo que la
hipótesis \(|L(v)|\ge d^+(v)+1\) se preserva; los vértices fuera de
\(S\) conservan sus listas y sus grados de salida no aumentan.
Repetimos. Todo vértice no coloreado mantiene siempre una lista no
vacía, y cada ronda con \(S\ne\varnothing\) colorea el conjunto no vacío
\(K\), por lo que finalmente todos los vértices quedan coloreados.
\(\square\)

\subsection{D.2 Emparejamientos
estables}\label{d.2-emparejamientos-estables}

Sea \(B\) un grafo bipartito con partes \(U\) y \(R\) en el que cada
vértice \(z\) tiene un orden lineal de preferencia \(>_z\) sobre sus
aristas incidentes. Un matching \(M\subseteq E(B)\) es \textbf{estable}
si toda arista \(f\notin M\) tiene un extremo \(z\) cubierto por una
arista \(e\in M\) con \(e>_zf\).

\subsubsection{Lema D.2}\label{lema-d.2-galeshapley-13}

Para todo sistema de preferencias existe un matching estable.

\subsubsection{Demostración}\label{demostraciuxf3n-8}

Ejecutamos aceptación diferida. Mientras exista algún \(u\in U\) no
emparejado que todavía no haya propuesto a lo largo de todas sus
aristas, hacemos que \(u\) proponga por su arista no intentada más
preferida, digamos \(f\), con extremo \(r\in R\). Si \(r\) no está
emparejado, o prefiere \(f\) a la arista que mantiene actualmente,
entonces \(r\) acepta \(f\), liberando a su pareja anterior, si existe;
en caso contrario, \(r\) rechaza \(f\). Cada arista recibe a lo sumo una
propuesta, de modo que el proceso termina, y termina en un matching
\(M\).

Sea \(f=ur\notin M\). Supongamos primero que \(u\) propuso a lo largo de
\(f\) en algún momento. Entonces \(r\) rechazó \(f\) inmediatamente o la
aceptó y más tarde la liberó; en ambos casos, en ese momento \(r\)
mantenía o recibió una arista que prefiere estrictamente a \(f\). Como
la arista mantenida por \(r\) solo mejora durante el proceso, la arista
final \(e\in M\) en \(r\) satisface \(e>_rf\). Supongamos ahora que
\(u\) nunca propuso a lo largo de \(f\). Un vértice no emparejado de
\(U\) sigue proponiendo mientras queden aristas no intentadas, por lo
que al terminar \(u\) está emparejado, digamos con \(e\in M\); como
\(u\) propone en orden decreciente de preferencia y nunca llegó a \(f\),
se tiene \(e>_uf\). En ambos casos, \(f\) tiene un extremo cuya arista
emparejada se prefiere a \(f\). \(\square\)

\subsection{D.3 El teorema de coloreo}\label{d.3-el-teorema-de-coloreo}

\subsubsection{Teorema D.3}\label{teorema-d.3-galvin-caso-de-grado-muxe1ximo}

Sea \(B\) un grafo bipartito simple de grado máximo \(\Delta\), y
supongamos que cada arista \(e\) lleva una lista \(L(e)\) con
\(|L(e)|\ge\Delta\). Entonces \(B\) admite un coloreo propio de aristas
que elige cada color desde la lista correspondiente.

\subsubsection{Demostración}\label{demostraciuxf3n-9}

\textbf{Paso 1: un coloreo propio de aristas con \(\Delta\) colores.}
Por el teorema de coloreo de aristas de König {[}14{]}, \(B\) admite un
coloreo propio de aristas \(\varphi:E(B)\to\{1,\ldots,\Delta\}\). Para
completar el argumento, procedemos por inducción en el número de
aristas. Eliminamos una arista \(e=ur\) y coloreamos el resto. A lo sumo
\(\Delta-1\) colores aparecen en \(u\) y a lo sumo \(\Delta-1\) en
\(r\), de modo que a \(u\) le falta un color \(\alpha\) y a \(r\) un
color \(\beta\). Si les falta un color común, se lo asignamos a \(e\).
En caso contrario, consideramos el camino maximal \(P\) que parte de
\(u\) y cuyas aristas están coloreadas alternadamente \(\beta,\alpha\).
El camino \(P\) no puede terminar en \(r\): como a \(r\) le falta
\(\beta\), tal camino terminaría con una arista de color \(\alpha\) y
tendría longitud par, mientras todo camino entre los vértices adyacentes
\(u\) y \(r\) de un grafo bipartito tiene longitud impar. Intercambiamos
los colores \(\alpha\) y \(\beta\) a lo largo de \(P\); por maximalidad
de \(P\), el coloreo sigue siendo propio, \(r\) no se altera y ahora a
\(u\) le falta \(\beta\). Asignamos \(\beta\) a \(e\).

\textbf{Paso 2: la orientación.} Definimos preferencias a partir de
\(\varphi\): cada \(u\in U\) prefiere sus aristas incidentes con color
\(\varphi\) \textbf{mayor}, y cada \(r\in R\) prefiere color \(\varphi\)
\textbf{menor}. Definimos un dígrafo \(D\) sobre el conjunto de vértices
\(E(B)\): para aristas distintas \(e,f\) que comparten un extremo \(z\),
ponemos un arco \(e\to f\) exactamente cuando \(f>_ze\).

Grados de salida: sea \(e=ur\) con \(\varphi(e)=c\). Los vecinos de
salida de \(e\) son las aristas en \(u\) con color mayor que \(c\) ---a
lo sumo \(\Delta-c\), pues los colores en \(u\) son dos a dos
distintos--- y las aristas en \(r\) con color menor que \(c\) ---a lo
sumo \(c-1\). Por tanto,

\[
d^+_D(e)\le(\Delta-c)+(c-1)=\Delta-1.
\]

\textbf{Paso 3: perfección por núcleos.} Un subdígrafo inducido de \(D\)
es \(D[S]\) para un conjunto \(S\) de aristas de \(B\). Un núcleo de
\(D[S]\) es exactamente un matching estable del subgrafo \((V(B),S)\)
con las preferencias heredadas: independencia en \(D\) significa que dos
aristas del núcleo no comparten extremo, pues dos aristas en un mismo
vértice siempre están unidas por un arco; y la condición de dominación
pide que todo \(f\in S\setminus K\) tenga un arco hacia \(K\), es decir,
un extremo \(z\) y una arista \(e\in K\) en \(z\) con \(e>_zf\), que es
precisamente estabilidad. El Lema D.2 proporciona un matching estable,
de modo que \(D\) es perfecto por núcleos.

\textbf{Paso 4: conclusión.} Toda arista satisface
\(|L(e)|\ge\Delta\ge d^+_D(e)+1\), de modo que el Lema D.1 da un coloreo
propio del grafo subyacente de \(D\) desde las listas; esto es, un
coloreo propio de aristas por listas de \(B\). \(\square\)

\subsubsection{Observación D.4 ---
Aplicación}\label{observaciuxf3n-d.4-aplicaciuxf3n}

El grafo de ganancias de la Sección 7.2 es simple y bipartito, y la
hipótesis (7.2) da para toda lista

\[
b-t_i\ge\max\{\rho,u\}\ge\Delta
\]

del grafo de ganancias. Por tanto, el Teorema D.3 se aplica
literalmente. La versión para multigrafos del teorema de Galvin también
es válida {[}12{]}, pero no se necesita aquí.

\begin{center}\rule{0.5\linewidth}{0.5pt}\end{center}

\subsection{Agradecimientos}\label{agradecimientos}

El autor agradece profundamente a su esposa María Paz y a sus hijos
Lucas, Juan Cristóbal, Francisca, Raimundo y Benjamín por su amor,
paciencia y apoyo.

\begin{center}\rule{0.5\linewidth}{0.5pt}\end{center}

\subsection{Uso de herramientas asistidas por
IA}\label{uso-de-herramientas-asistidas-por-ia}

Se utilizaron herramientas asistidas por IA durante las etapas
exploratoria, computacional, adversarial, organizativa y editorial,
incluidos sistemas de Anthropic, Google y OpenAI. Estas herramientas
apoyaron la prueba de argumentos candidatos, las comprobaciones exactas
de regresión, la organización de la demostración, la preparación de
auditorías y la redacción. El autor revisó el contenido matemático,
seleccionó los argumentos finales y conserva la responsabilidad
exclusiva por las afirmaciones, citas, código y presentación. Ningún
sistema de IA figura como autor.

\begin{center}\rule{0.5\linewidth}{0.5pt}\end{center}

\subsection{Referencias}\label{referencias}

{[}1{]} B. Barber, D. Kühn, A. Lo, and D. Osthus, ``Edge-decompositions
of graphs with high minimum degree,'' \emph{Advances in Mathematics}
\textbf{288} (2016), 337--385.

{[}2{]} O. V. Borodin, A. V. Kostochka, and D. R. Woodall, ``List edge
and list total colourings of multigraphs,'' \emph{Journal of
Combinatorial Theory, Series B} \textbf{71} (1997), no. 2, 184--204.

{[}3{]} G.-T. Chen, P. Erdős, and E. T. Ordman, ``Clique partitions of
split graphs,'' in \emph{Combinatorics, Graph Theory, Algorithms and
Applications} (Beijing, 1993), World Scientific, 1994, pp.~21--30.

{[}4{]} G. A. Dirac, ``Some theorems on abstract graphs,''
\emph{Proceedings of the London Mathematical Society} (3) \textbf{2}
(1952), 69--81.

{[}5{]} F. Dross, ``Fractional triangle decompositions in graphs with
large minimum degree,'' \emph{SIAM Journal on Discrete Mathematics}
\textbf{30} (2016), no. 1, 36--42.

{[}6{]} P. Erdős, E. T. Ordman, and Y. Zalcstein, ``Clique partitions of
chordal graphs,'' \emph{Combinatorics, Probability and Computing}
\textbf{2} (1993), no. 4, 409--415.

{[}7{]} P. E. Haxell and V. Rödl, ``Integer and fractional packings in
dense graphs,'' \emph{Combinatorica} \textbf{21} (2001), 13--38.

{[}8{]} J. P. Traverso Gianini, ``Complete-split extremizers for a
fractional triangle-cover functional on chordal graphs'' (Paper II in
the series), preprint v1.0, July 2026.

{[}9{]} R. Yuster, ``Integer and fractional packing of families of
graphs,'' \emph{Random Structures \& Algorithms} \textbf{26} (2005),
110--118.

{[}10{]} P. Keevash, ``The existence of designs,'' arXiv:1401.3665,
revised 2018.

{[}11{]} P. Allen, J. Böttcher, O. Cooley, and R. Mycroft, ``Tight
cycles and regular slices in dense hypergraphs,'' \emph{Journal of
Combinatorial Theory, Series A} \textbf{149} (2017), 30--100.

{[}12{]} F. Galvin, ``The list chromatic index of a bipartite
multigraph,'' \emph{Journal of Combinatorial Theory, Series B}
\textbf{63} (1995), no. 1, 153--158.

{[}13{]} D. Gale and L. S. Shapley, ``College admissions and the
stability of marriage,'' \emph{The American Mathematical Monthly}
\textbf{69} (1962), no. 1, 9--15.

{[}14{]} D. König, ``Über Graphen und ihre Anwendung auf
Determinantentheorie und Mengenlehre,'' \emph{Mathematische Annalen}
\textbf{77} (1916), 453--465.

{[}15{]} J. P. Traverso Gianini, ``Affine profile reduction for
fractional triangle packings in split graphs'' (Paper I in the series),
preprint v1.0, July 2026.

\end{document}
